\section{The non-normal fibre and the CDP obstruction}

Campana--Demailly--Peternell prove that a compact complex threefold with
$H^1(X;\mathbb Z)=H^2(X;\mathbb Z)=0$, algebraic dimension one, and
holomorphic algebraic reduction has $c_3(X)\le0$
\cite[Theorem~2.2, p.~680]{CDP2020}.  Their corollary gives $a(X)=0$ for
a compact complex threefold homeomorphic to $S^6$
\cite[Corollary~2.3, p.~680]{CDP2020}.  The candidate asserts all those
hypotheses together with $c_3(X)=2$ and $a(X)=1$.  This section separates:
\begin{enumerate}[label=(\roman*)]
\item validity of one printed inference on p.~692 of \cite{CDP2020};
\item validity of the published theorem by any other argument;
\item validity of the candidate's proposed non-normal-fibre mechanism; and
\item truth of the candidate's upstream construction and topology.
\end{enumerate}

\subsection{The exact p.~692 inference}

Let $S$ be the reduction of a fibre in a smooth threefold $\mathcal X$,
and let $A$ denote the line-bundle twist after duality.  The two relevant
sequences are
\[
0\longrightarrow N^*_{S/\mathcal X}\otimes A
\longrightarrow\Omega^1_{\mathcal X}|_S\otimes A
\xrightarrow{\rho}\Omega^1_S\otimes A\longrightarrow0
\tag{5.1}
\]
and
\[
0\longrightarrow\mathcal T(\Omega^1_S)\otimes A
\longrightarrow\Omega^1_S\otimes A
\longrightarrow\widetilde\Omega^1_S\otimes A\longrightarrow0,
\tag{5.2}
\]
where $\mathcal T$ denotes sections supported on a nowhere dense analytic
set.  The opening reduction of \cite[\S7, p.~692]{CDP2020} treats the
vanishing of the global sections of the left end of (5.1) and the right end
of (5.2) as sufficient for the vanishing of the middle term of (5.1).
That inference leaves untreated a section whose image under $\rho$ is a
nonzero torsion section.  The omitted sheaf can be identified exactly.  Put
\[
 \mathscr T_S:=\ker\bigl(\Omega_S^1\to
 j_*\Omega^1_{S_{\rm reg}}\bigr),
 \qquad \widetilde\Omega_S^1=\Omega_S^1/\mathscr T_S,
\]
and
\[
 \mathscr K_S:=
 \ker\bigl(\Omega_X^1|_S\longrightarrow\widetilde\Omega_S^1\bigr).
\]

\begin{proposition}[the missing cotangent filtration]
\label{prop:missing-cotangent-filtration}
Let $S$ be a reduced effective Cartier divisor in a smooth complex
threefold $X$.  There are natural short exact sequences
\[
0\longrightarrow N^*_{S/X}\longrightarrow\mathscr K_S
\longrightarrow\mathscr T_S\longrightarrow0,
\]
\[
0\longrightarrow\mathscr K_S\longrightarrow\Omega_X^1|_S
\longrightarrow\widetilde\Omega_S^1\longrightarrow0.
\]
\end{proposition}

\begin{proof}
The conormal sequence is right exact.  It is also left exact here.  Locally
write $S=(g=0)$ with $g$ reduced.  If $h\,dg|_S=0$, then $h$ vanishes at the
generic smooth point of every irreducible component, because $dg$ is nonzero
there.  The ring $\mathcal O_S$ is reduced, so the intersection of its
minimal primes is zero; hence $h=0$.  Thus
\[
0\to N^*_{S/X}\to\Omega_X^1|_S\to\Omega_S^1\to0
\]
is exact.  By definition, $\mathscr K_S$ is the inverse image of
$\mathscr T_S$ under the last surjection.  Restriction to that inverse image
is onto $\mathscr T_S$ and has kernel $N^*_{S/X}$, proving the first
sequence.  Composing with the quotient by $\mathscr T_S$ proves the second.
\end{proof}

\begin{proposition}[the exact global-section criterion]
\label{prop:missing-global-section-criterion}
Let $A$ be a line bundle on $S$.  If
\[
H^0(S,\widetilde\Omega_S^1\otimes A)=0,
\]
then
\[
H^0(S,\Omega_X^1|_S\otimes A)
=H^0(S,\mathscr K_S\otimes A).
\]
If also $H^0(S,N^*_{S/X}\otimes A)=0$, then this group is canonically
\[
\ker\!\left[
H^0(S,\mathscr T_S\otimes A)
\xrightarrow{\delta_A}H^1(S,N^*_{S/X}\otimes A)
\right].
\]
Hence the two vanishings printed in the p.~692 reduction imply the desired
ambient vanishing if and only if $\delta_A$ is injective.
\end{proposition}

\begin{proof}
Tensoring the second sequence of
\cref{prop:missing-cotangent-filtration} by the locally free sheaf $A$
preserves exactness; its global-section sequence gives the first equality.
Tensor the first sequence and take global sections.  If the conormal
$H^0$ vanishes, exactness identifies the middle group with the displayed
kernel.  No condition on the torsion-free quotient controls that kernel.
\end{proof}

\subsection{Local algebra of a reduced normal-crossings fibre}

Let $g:\mathcal X\to\Delta$ be a proper holomorphic map from a smooth complex
threefold to a disc, and put $S=g^{-1}(0)$.  Assume $S$ is reduced and has
normal crossings: at every point there are coordinates $(x,y,z)$ in which
\[
g=ux,\qquad g=uxy,\qquad\text{or}\qquad g=uxyz
\tag{5.3}
\]
for a unit $u$.  Let $\eta:\widetilde S\to S$ be the normalization,
$D=S_{\mathrm{sing,red}}$, and $\widetilde D=\eta^{-1}(D)_{\mathrm{red}}$.

\begin{lemma}[the conormal map is injective]\label{lem:conormal-injective}
\Verified\ In the setting above, $N^*_{S/\mathcal X}\simeq\mathcal O_S$,
and its generator maps to $dg|_S$.  Multiplication by $dg|_S$ is injective,
so
\[
0\longrightarrow\mathcal O_S
\xrightarrow{\,h\mapsto h\,dg\,}
\Omega^1_{\mathcal X}|_S\longrightarrow\Omega^1_S\longrightarrow0
\tag{5.4}
\]
is exact.
\end{lemma}

\begin{proof}
The fibre ideal is the principal ideal $(g)$ and is reduced, so
$(g)/(g^2)\simeq\mathcal O_S$.  Multiplication by a unit does not affect
injectivity, and the term $f\,du$ vanishes after restriction to $f=0$.
It therefore suffices to use $f=x,xy,xyz$.

For $f=xy$, if a lift $\widetilde h$ satisfies
$\widetilde h(y\,dx+x\,dy)=0$ modulo $(xy)$, then
$\widetilde h y,\widetilde h x\in(xy)$.  Hence
$\widetilde h\in(x)\cap(y)=(xy)$ in the regular local ring, so $h=0$ on
$S$.  For $f=xyz$, the three coefficient conditions say
$\widetilde h yz,\widetilde h xz,\widetilde h xy\in(xyz)$; hence
$\widetilde h\in(x)\cap(y)\cap(z)=(xyz)$.  The case $f=x$ is immediate.
Thus the conormal map is injective in every local model, proving (5.4).
\end{proof}

\begin{lemma}[normalization--conductor exact sequence]
\label{lem:conductor-square}
\Verified\ The sequence
\[
0\longrightarrow\mathcal O_S
\xrightarrow{\alpha}
\eta_*\mathcal O_{\widetilde S}\oplus\mathcal O_D
\xrightarrow{\beta}
\eta_*\mathcal O_{\widetilde D}
\longrightarrow0,
\tag{5.5}
\]
\[
\alpha(h)=(\eta^*h,h|_D),\qquad
\beta(\widetilde h,k)=
\widetilde h|_{\widetilde D}-\eta^*k,
\]
is exact.
\end{lemma}

\begin{proof}
The assertion is local.  Away from $D$, $\eta$ is an isomorphism.
At a double point $S=\{xy=0\}$, a member of $\ker\beta$ is
\[
\bigl((h_1(y,z),h_2(x,z)),k(z)\bigr),
\qquad
h_1(0,z)=h_2(0,z)=k(z).
\]
The function
\[
h=h_1(y,z)+h_2(x,z)-h_1(0,z)
\]
on $\{xy=0\}$ restricts to $h_1,h_2$ on the two branches and to $k$ on
their intersection.  This proves middle exactness in the double model.

At a triple point $S=\{xyz=0\}$, write the branch functions as
$h_1(y,z)$, $h_2(x,z)$, $h_3(x,y)$.  Membership in $\ker\beta$ is exactly
\[
h_1(0,z)=h_2(0,z),\quad
h_1(y,0)=h_3(0,y),\quad
h_2(x,0)=h_3(x,0).
\]
Then
\[
\begin{split}
h={}&h_1(y,z)+h_2(x,z)+h_3(x,y)\\
&-h_1(0,z)-h_1(y,0)-h_2(x,0)+h_1(0,0)
\end{split}
\]
restricts to the prescribed function on every branch.  Injectivity of
$\alpha$ follows because $S$ is reduced and $\eta$ is an isomorphism over a
dense open set.  Surjectivity of $\beta$ follows locally by lifting a
function on the branch conductor to the smooth branch and taking $k=0$.
This proves exactness in every local model.
\end{proof}

\begin{theorem}[conductor-supported ambient differential]
\label{prop:conductor}
\Verified\ In the setting (5.3), let $A$ be a line bundle on $S$ with
$\eta^*A\simeq\mathcal O_{\widetilde S}$, and choose a nowhere-vanishing
section $e$ of $\eta^*A$.  There is a unique nonzero
\[
s\in H^0(S,\Omega^1_{\mathcal X}|_S\otimes A)
\tag{5.6}
\]
such that
\[
\eta^*s=\eta^*(dg|_S)\otimes e,\qquad s|_D=0.
\tag{5.7}
\]
Its image $\rho(s)$ in $\Omega^1_S\otimes A$ is supported on $D$ and
therefore torsion; its image in
$\widetilde\Omega^1_S\otimes A$ is zero.  Moreover:
\begin{enumerate}[label=(\alph*)]
\item $\rho(s)=0$ if and only if $e=\eta^*\sigma$ for a nowhere-vanishing
$\sigma\in H^0(S,A)$;
\item if $S$ is connected and $A\not\simeq\mathcal O_S$, then
$H^0(S,A)=0$ and $s$ does not come from the conormal term in (5.1).
\end{enumerate}
\end{theorem}

\begin{proof}
In the double model, $d(xy)=y\,dx+x\,dy$ has coefficients in the reduced
ideal $(x,y)$ of $D$.  In the triple model,
\[
d(xyz)=yz\,dx+xz\,dy+xy\,dz
\]
has coefficients in $(xy,xz,yz)$, the reduced ideal of the union of the
three double curves.  Unit factors in (5.3) do not change this after
restriction to $S$.  Hence $\eta^*(dg|_S)$ restricts to zero on
$\widetilde D$.  Tensor (5.5) by the locally free sheaf
$E=\Omega^1_{\mathcal X}|_S\otimes A$.  The compatible pair
\[
\bigl(\eta^*(dg|_S)\otimes e,0\bigr)
\]
therefore descends uniquely to $s$.  On the smooth locus of $S$, $dg|_S$
is a nonzero conormal vector and $e$ is nowhere zero, so $s\ne0$.

On $S\setminus D$, the normalization is an isomorphism and
$s=dg|_S\otimes e$.  By (5.4) its intrinsic image vanishes there.  Thus
$\rho(s)$ is supported on $D$ and lies in the torsion subsheaf; its image
modulo torsion is zero.

If $\rho(s)=0$, exactness of (5.4) after tensoring by $A$ gives
$s=dg|_S\otimes\sigma$ for a section $\sigma$ of $A$.  Pulling back and
comparing with (5.7) on the dense open set where $dg\ne0$ gives
$\eta^*\sigma=e$ there, hence everywhere by holomorphic continuation.
Since $e$ never vanishes and $\eta$ is surjective, $\sigma$ never vanishes
and trivializes $A$.  The converse is immediate.

Finally suppose $S$ is connected and $\sigma\in H^0(S,A)$.  On each compact
connected component of $\widetilde S$, write
$\eta^*\sigma=c_i e$; the holomorphic function $c_i$ is constant.  If one
$c_i$ is zero and another is nonzero, connectedness of the finite union of
irreducible components of $S$ provides two such components meeting at a
point.  The value of $\sigma$ at that point would be both zero and nonzero,
a contradiction.  Hence either every $c_i=0$, so $\sigma=0$, or every
$c_i\ne0$, so $\sigma$ is nowhere zero and $A$ is trivial.  When $A$ is
nontrivial only the first case occurs.  Since
$N^*_{S/\mathcal X}\simeq\mathcal O_S$, the conormal global-section space is
$H^0(S,A)=0$, while $s\ne0$.
\end{proof}

\subsection{A proper counterexample to the printed inference}

\begin{theorem}[nodal-cubic test family]\label{prop:nodal-counterexample}
\Refuted\ The p.~692 implication described after (5.2) is false as a formal
sheaf-theoretic statement, even for a proper family with smooth total space
and an irreducible reduced normal-crossings central fibre.
\end{theorem}

\begin{proof}
Let $\Delta$ be a sufficiently small disc with coordinate $t$ and put
\[
\mathcal C=
\{([x:y:z],t)\in\mathbb P^2\times\Delta:
y^2z=x^2(x+z)+tz^3\}.
\tag{5.8}
\]
This is proper over $\Delta$.  Its total space is smooth.  On $z\ne0$, the
partial derivative of the defining equation with respect to $t$ is
$-z^3\ne0$.  On $z=0$, the equation forces $x=0$; then $y\ne0$ and the
$z$-derivative is $y^2\ne0$.  The central fibre $C_0$ is the irreducible
nodal cubic $y^2z=x^2(x+z)$, with node $[0:0:1]$.

Set
\[
\mathcal X=\mathcal C\times\mathbb P^1,\qquad
S=C_0\times\mathbb P^1,
\]
and let $g:\mathcal X\to\Delta$ be projection.  Then $\mathcal X$ is a
smooth threefold, $g$ is proper, and $S$ is an irreducible reduced
normal-crossings surface.  Its normalization is
$\widetilde S\simeq\mathbb P^1\times\mathbb P^1$.

Let $p,q$ be the two inverse images of the node in the normalization of
$C_0$.  Fix $\lambda\in\mathbb C^*$ which is not a root of unity (for
example, $\lambda=2$).  Glue the fibres of
the trivial line bundle on the normalization at $p$ and $q$ by the scalar
$\lambda$.  This defines a line bundle
$A_\lambda\in\operatorname{Pic}^0(C_0)$ whose pullback is trivial.  It has
infinite order: its $n$th tensor power has gluing scalar $\lambda^n$, which
is never one.  It is not trivial: a global section pulls back to a constant
$c$ satisfying
$c=\lambda c$, so
\[
H^0(C_0,A_\lambda)=0.
\tag{5.9}
\]
Put $A=\operatorname{pr}_1^*A_\lambda$ on $S$.  Then
$\eta^*A\simeq\mathcal O_{\widetilde S}$ and $H^0(S,A)=0$.
Because $S$ is the full fibre,
$N^*_{S/\mathcal X}\simeq\mathcal O_S$, so the conormal-twist global
sections vanish.

The natural map
$\widetilde\Omega^1_S\to\eta_*\Omega^1_{\widetilde S}$ is injective:
its kernel is supported where normalization is not an isomorphism, while
the source is torsion-free by definition.  Projection formula and
$\eta^*A\simeq\mathcal O_{\widetilde S}$ therefore give
\[
H^0(S,\widetilde\Omega^1_S\otimes A)
\hookrightarrow
H^0(\widetilde S,\Omega^1_{\widetilde S})=0,
\tag{5.10}
\]
the last equality following from
$\widetilde S=\mathbb P^1\times\mathbb P^1$.
Nevertheless \cref{prop:conductor} produces
\[
0\ne s\in H^0(S,\Omega^1_{\mathcal X}|_S\otimes A).
\]
Thus both vanishings invoked in the printed reduction hold while the desired
ambient vanishing fails.  This test chooses the line bundle $A$ directly on
$S$; it does not assert that $A$ extends to a line bundle on the total space
$\mathcal X$.  It therefore refutes the claimed \emph{formal implication
between the three sheaf-cohomology groups} for a non-torsion $A$, but does
not by itself instantiate the narrower global-line-bundle situation on
p.~692.  The family (5.8) also does not satisfy all hypotheses of CDP
Theorem~2.2, so this is not a counterexample to that theorem.
\end{proof}

\subsection{Cartier-divisor duality without a normality hypothesis}

\begin{theorem}[complete Cartier-fibre duality]
\label{thm:cartier-duality}
\Verified\ Let $X$ be a compact smooth complex threefold,
$i:W\hookrightarrow X$ an effective Cartier divisor, and $E$ a holomorphic
vector bundle.  Put $F=E|_W$ and
\[
\omega_W=(\omega_X\otimes\mathcal O_X(W))|_W.
\]
For $q=0,1,2$ there are natural perfect pairings
\[
H^q(W,F)^\vee
\simeq H^{2-q}(W,F^\vee\otimes\omega_W).
\tag{5.11}
\]
Moreover $H^q(W,F)=0$ for $q>2$.  No normality or reducedness hypothesis on
$W$ is used.
\end{theorem}

\begin{proof}
The Cartier divisor gives two short exact sequences
\[
0\to E(-W)\to E\to i_*F\to0
\tag{5.12}
\]
and, by adjunction,
\[
0\to E^\vee\otimes\omega_X
\to E^\vee\otimes\omega_X(W)
\to i_*(F^\vee\otimes\omega_W)\to0.
\tag{5.13}
\]
Smooth Serre duality on the compact threefold gives functorial isomorphisms
\[
H^r(X,G)^\vee\simeq H^{3-r}(X,G^\vee\otimes\omega_X)
\tag{5.14}
\]
for every vector bundle $G$
\cite[Theorem~9.37]{Lee2024}.  Naturality of the trace pairing says that the
dual of the map on cohomology induced by $E(-W)\to E$ is the map induced by
$E^\vee\otimes\omega_X\to E^\vee\otimes\omega_X(W)$.

Fix $q\in\{0,1,2\}$.  Dualize the relevant five terms of the long exact
sequence of (5.12) and use (5.14).  They become
\[
\begin{split}
H^{2-q}(X,E^\vee\otimes\omega_X)
&\longrightarrow H^{2-q}(X,E^\vee\otimes\omega_X(W))\\
&\longrightarrow H^q(W,F)^\vee\\
&\longrightarrow H^{3-q}(X,E^\vee\otimes\omega_X)
\longrightarrow H^{3-q}(X,E^\vee\otimes\omega_X(W)).
\end{split}
\tag{5.15}
\]
The long exact sequence of (5.13) is the same five-term sequence with its
middle group replaced by
$H^{2-q}(W,F^\vee\otimes\omega_W)$.  Naturality supplies a commutative
diagram whose four outer vertical arrows are identities.  The five lemma
(or exactness at the three middle terms) gives the unique middle isomorphism
(5.11).  Finite dimensionality makes the pairing perfect.

For $q=3$, the map
$H^3(X,E(-W))\to H^3(X,E)$ is dual under (5.14) to the injective map
\[
H^0(X,E^\vee\otimes\omega_X)
\hookrightarrow H^0(X,E^\vee\otimes\omega_X(W)).
\]
It is therefore surjective.  The long exact sequence of (5.12) gives
$H^3(W,F)=0$.  For $q>3$ both adjacent ambient cohomology groups vanish by
the dimension-three Dolbeault resolution, so (5.12) gives
$H^q(W,F)=0$.  This also agrees with analytic duality for the Gorenstein
hypersurface $W$ in Ramis--Ruget
\cite[Theorems~1--2 and Proposition~2]{RamisRuget1970}.
\end{proof}

\begin{proposition}[absolute duality on a compact Cartier fibre]
\label{prop:cartier-duality-absolute}
\Verified\ Let $X$ be a smooth complex threefold, not necessarily compact,
and let $i:W\hookrightarrow X$ be a compact effective Cartier divisor that
is the full scheme-theoretic fibre of a holomorphic map $f:X\to B$ to a
smooth curve.  For every vector bundle $E$ on $X$ there is a functorial
perfect pairing
\[
 H^2(W,E|_W)\times
 H^0(W,E^\vee|_W\otimes\omega_X|_W)\longrightarrow\mathbb C.
\]
Equivalently,
\[
 H^2(W,E|_W)^\vee\simeq
 H^0(W,E^\vee|_W\otimes\omega_X|_W).
\tag{5.15a}
\]
Neither normality of $W$ nor compactness of the ambient threefold is
required.
\end{proposition}

\begin{proof}
An effective Cartier divisor in a smooth threefold is locally a
two-dimensional hypersurface quotient of a regular local ring.  Hence $W$
is Cohen--Macaulay and Gorenstein.  If $K_W^\bullet$ denotes its analytic
dualizing complex, the complete-intersection calculation of Ramis--Ruget
gives
\[
 K_W^\bullet\simeq\omega_W[2]
\]
with $\omega_W$ invertible
\cite[Proposition~2 and the complete-intersection calculation,
pp.~90--91]{RamisRuget1970}.  Cartier adjunction gives
\[
 \omega_W\simeq
 (\omega_X\otimes\mathcal O_X(W))|_W.
\]
Because $W$ is the full fibre over $b\in B$, one has
$\mathcal O_X(W)\simeq f^*\mathcal O_B(b)$ locally along $W$; its
restriction to $W$ is the pullback of a one-dimensional vector space.
Choosing a nonzero vector identifies
$\mathcal O_X(W)|_W\simeq\mathcal O_W$ and
$\omega_W\simeq\omega_X|_W$.  A different choice rescales the final trace
pairing by a nonzero scalar.

Ramis--Ruget duality pairs $H_c^p(W,\mathcal G)$ with
$\operatorname{Ext}^{-p}_W(\mathcal G,K_W^\bullet)$
\cite[Theorem~2, pp.~82--83]{RamisRuget1970}.  Since $W$ is compact,
$H_c^p=H^p$.  Taking $p=2$ and $\mathcal G=E|_W$, the shift and local
freeness give
\[
 \operatorname{Ext}^{-2}_W(\mathcal G,K_W^\bullet)
 \simeq\operatorname{Hom}_W(\mathcal G,\omega_W)
 =H^0(W,E^\vee|_W\otimes\omega_W).
\]
Grauert finiteness makes both coherent cohomology spaces finite-dimensional
\cite{Grauert1960}; the strong analytic duality is therefore the ordinary
algebraic perfect pairing (5.15a).  Its variance is contravariant in $E|_W$
by precomposition on the Hom term.  This proves the assertion.
\end{proof}

\subsection{Flatness and top-degree analytic base change}

\begin{lemma}[automatic properness, flatness, and Cartier fibres]
\label{lem:curve-flatness}
\Verified\ Let $f:X\to B$ be a nonconstant holomorphic map from a connected
complex manifold to a Riemann surface.  Then $f$ is flat, every holomorphic
vector bundle on $X$ is $f$-flat, and every full scheme-theoretic fibre is an
effective Cartier divisor of pure codimension one.  If $X$ is compact, then
$f$ is proper.
\end{lemma}

\begin{proof}
Fix $x\in X$, $c=f(x)$, and put
$R=\mathcal O_{B,c}$, $S=\mathcal O_{X,x}$.  The ring $R$ is a discrete
valuation ring and $S$ is a regular local domain.  The homomorphism
$R\to S$ is injective.  Indeed, if a nonzero germ $h\in R$ mapped to zero,
$h\circ f$ would vanish on a neighbourhood of $x$.  The zero set of a
nonzero one-variable holomorphic germ is isolated, so $f$ would be constant
on that neighbourhood and therefore, by the identity theorem on connected
$X$, constant everywhere.

Thus every nonzero element of $R$ acts on the domain $S$ as a nonzero
element and hence as a nonzerodivisor.  Therefore $S$ is torsion-free as an
$R$-module.  A module over a discrete valuation ring is flat if and only if
it is torsion-free \cite[Tag~0AUW]{StacksProject}; hence $S$ is $R$-flat.
This holds at every point, so $f$ is flat.  A vector bundle is locally a
finite free $\mathcal O_X$-module, and a finite direct sum of $R$-flat
modules is flat; hence it is $f$-flat.  If $t$ is a local parameter at
$c$, then $f^*t$ is a nonzero element of the domain $S$ and thus a
nonzerodivisor.  The full scheme-theoretic fibre is therefore locally cut
out by one nonzerodivisor, so it is an effective Cartier divisor of pure
codimension one.  Finally, when $X$ is compact, the underlying continuous
map to the Hausdorff surface $B$ is proper.
\end{proof}

\begin{proposition}[the exact analytic base-change input]
\label{prop:analytic-base-change}
\Verified\ Let $f:X\to B$ be proper between complex spaces, let
$\mathcal F$ be coherent and $f$-flat, and fix $c\in B$.  If
\[
H^3(X_c,\mathcal F|_{X_c})=0,
\tag{5.16}
\]
then the canonical fibre map
\[
\beta_c^2:
(R^2f_*\mathcal F)_c\otimes_{\mathcal O_{B,c}}\mathbb C(c)
\longrightarrow H^2(X_c,\mathcal F|_{X_c})
\tag{5.17}
\]
is an isomorphism.
\end{proposition}

\begin{proof}
Apply B\u{a}nic\u{a}--St\u{a}n\u{a}\c{s}il\u{a},
Chapter~III, \S3, Corollary~3.5, with $q=3$
\cite[pp.~120--121]{BanicaStanasila1976}.  Its hypotheses are precisely
properness, coherence, and flatness.  Condition (5.16) first forces the
degree-three direct image to vanish near $c$.  The corollary then applies the
equivalent base-change conditions of its preceding Theorem~3.4 in degree
$q-1=2$ and states that the canonical map (5.17) is bijective.  This is the
top-degree argument used in \cite[Lemma~1.5]{CDP1998}; it does not require
the stronger, false assertion that all direct images are locally free.
\end{proof}

\begin{corollary}[uniform top-degree base change for surface fibres]
\label{cor:surface-top-base-change}
\Verified\ Let $f:X\to B$ be a proper flat holomorphic map to a Riemann
surface whose full scheme-theoretic fibres have pure complex dimension two,
and let $E$ be a holomorphic vector bundle on $X$.  Then
\[
 R^qf_*E=0\quad(q>2)
 \tag{5.17-V}
\]
and, for every $b\in B$, including a point with nonreduced fibre, the
canonical map
\[
 (R^2f_*E)_b\otimes_{\mathcal O_{B,b}}\mathbb C(b)
 \xrightarrow{\ \sim\ }H^2(X_b,E|_{X_b})
 \tag{5.17-T}
\]
is an isomorphism.  No constancy of $h^2$, local freeness of $R^2f_*E$, or
reducedness of $X_b$ is required.
\end{corollary}

\begin{proof}
On a sufficiently small neighbourhood $U$ of $b$, the finite-complex theorem
used in analytic cohomology and base change represents $Rf_*E|_U$ by a
bounded complex $K^\bullet$ of finite free $\mathcal O_U$-modules
\cite[Chapter~III, \S3, Theorem~3.4]{BanicaStanasila1976}.  Flatness of
$E$ over $B$ identifies its derived fibre with the ordinary restriction to
the full scheme-theoretic fibre.  Over the discrete valuation ring
$A=\mathcal O_{B,b}$, the universal-coefficient sequence is
\[
0\longrightarrow H^q(K^\bullet)_b\otimes_A\mathbb C(b)
 \longrightarrow H^q(K^\bullet\otimes_A\mathbb C(b))
 \longrightarrow
 \operatorname{Tor}_1^A(H^{q+1}(K^\bullet)_b,\mathbb C(b))
 \longrightarrow0.
\tag{5.17-U}
\]
If some $R^qf_*E$ with $q>2$ were nonzero, choose the largest such $q$ and a
point of its support.  Nakayama's lemma makes the left term of (5.17-U)
nonzero there, while maximality kills the Tor term.  This would give nonzero
coherent cohomology in degree $q>2$ on a compact two-dimensional complex
space, a contradiction.  Thus (5.17-V) holds.  For $q=2$, the only
obstruction in (5.17-U) is
$\operatorname{Tor}_1^A((R^3f_*E)_b,\mathbb C(b))$, which is zero by
(5.17-V), proving (5.17-T).  This is the explicit obstruction calculation
behind \cref{prop:analytic-base-change}; it was independently supplied, with
eight named local certificates, in Reviewer~B's candidate-specific follow-up
\cite{ReviewerBValidation2026}.
\end{proof}

\subsection{A topology-independent \texorpdfstring{$\operatorname{Pic}^0$}{Pic0} route and three all-Pic routes}

Let $i:W\hookrightarrow X$ denote the cusp-fibre inclusion and
$\eta:\widetilde W\simeq\mathrm{dP}_6\to W$ its normalization.

\begin{proposition}[topologically trivial twists pull back trivially]
\label{prop:pic0-pullback}
\Verified\ Assume that $\omega_X$ is trivial on a neighbourhood of $W$.
For every $L\in\operatorname{Pic}^0(X)$, put
\[
 A_L=(L^\vee\otimes\omega_X)|_W.
\]
Then
\[
 \eta^*A_L\simeq\mathcal O_{\mathrm{dP}_6}.
\tag{5.17a}
\]
No assertion about $H^2(X;\mathbb Z)$ is used.
\end{proposition}

\begin{proof}
Every line bundle in $\operatorname{Pic}^0(X)$ has $c_1(L)=0$.  Naturality
therefore gives $c_1(\eta^*i^*L)=0$.  The degree-six del Pezzo surface is the
blow-up of $\mathbb P^2$ in three noncollinear points, so
$H^1(\mathrm{dP}_6,\mathcal O)=0$.  The exponential sequence consequently
makes
\[
 c_1:\operatorname{Pic}(\mathrm{dP}_6)\longrightarrow
 H^2(\mathrm{dP}_6;\mathbb Z)
\]
injective.  Hence $\eta^*i^*L$ is trivial.  The assumed neighbourhood
trivialization of $\omega_X$ also pulls back trivially, which proves
(5.17a).
\end{proof}

\begin{proposition}[the exact all-line-bundle restriction obstruction]
\label{prop:all-pic-restriction}
\Verified\ Define
\[
 r=\eta^*i^*:\operatorname{Pic}(X)\longrightarrow
 \operatorname{Pic}(\mathrm{dP}_6).
\]
Under the same local canonical trivialization,
\[
 \eta^*A_L\simeq r(L)^\vee
 \qquad(L\in\operatorname{Pic}(X)).
\tag{5.17b}
\]
Thus the trivial-pullback hypothesis used by
\cref{prop:conductor} holds simultaneously for every line bundle exactly
when $r=0$.  It always holds on $\operatorname{Pic}^0(X)$.
\end{proposition}

\begin{proof}
Pullback commutes with tensor product and duals, while the local
trivialization gives $\eta^*i^*\omega_X\simeq\mathcal O$.  This is (5.17b),
and its uniform equivalence with $r=0$ is immediate.  The final assertion is
\cref{prop:pic0-pullback}.
\end{proof}

\begin{proposition}[Leray-edge vanishing of normalized cusp restrictions]
\label{prop:leray-edge-restriction}
\Verified\ For the constructed map $f:X\to B$ and every
$L\in\operatorname{Pic}(X)$, the degree-two Leray edge section satisfies
\[
 e(c_1(L))=0\in H^0(B,R^2f_*\mathbb Z),
 \qquad
 \eta^*i^*L\simeq\mathcal O_{\mathrm{dP}_6}.
\tag{5.17c}
\]
This proof does not invoke the packaged conclusion
$H^2(X;\mathbb Z)=0$.
\end{proposition}

\begin{proof}
First record explicitly the compatibility used here.  For every $b\in B$,
proper base change and functoriality of the filtered-complex spectral
sequence give a commutative square
\[
\begin{CD}
 H^2(X;\mathbb Z) @>{e}>> H^0(B,R^2f_*\mathbb Z)\\
 @V{i_b^*}VV @VV{\operatorname{ev}_b}V\\
 H^2(X_b;\mathbb Z) @= H^2(X_b;\mathbb Z).
\end{CD}
\tag{5.17d}
\]
Indeed, pull the Leray filtered complex back along $\{b\}\to B$; the
spectral sequence over a point has only column zero, so its edge is the
identity, while the right vertical map is the proper-base-change
identification.  This uses the functoriality and convergence statements in
\cite[Tags~012O, 012U and~012W]{StacksProject} and analytic proper base
change \cite[Theorem~2.3.26]{Dimca2004}.  The connecting homomorphism in the
exponential sequence is natural under $i_b^*$, so (5.17d) sends
$c_1(L)$ to $c_1(L|_{X_b})$.

Put $s_L=e(c_1(L))$.  The image of an edge morphism consists of permanent
cycles, hence
\[
 s_L\in E_\infty^{0,2}\subseteq\ker
 \bigl(d_2^{0,2}:E_2^{0,2}\to E_2^{2,1}\bigr);
\tag{5.17e}
\]
see \cite[Tags~011N, 011O and~012U]{StacksProject}.  For the constructed
candidate, \cref{prop:leray-e2,lem:higher-transgressions} give
\[
 E_2^{0,2}=\mathbb Z(2q),\qquad
 E_2^{2,1}=\mathbb Z\,\omega[\delta],\qquad
 d_2^{0,2}(2q)=\omega[\delta].
\]
Thus $d_2^{0,2}$ is an isomorphism and (5.17e) forces $s_L=0$.
Taking $b=p_0$ in (5.17d) gives $c_1(i^*L)=0$, and hence
$c_1(\eta^*i^*L)=0$.  Finally, the injectivity of
$c_1:\operatorname{Pic}(\mathrm{dP}_6)\to H^2(\mathrm{dP}_6;\mathbb Z)$
proved in \cref{prop:pic0-pullback} makes $\eta^*i^*L$ trivial.
\end{proof}

\begin{lemma}[the pair-of-pants defect class]
\label{lem:picard-crossed-defect}
\Verified\ Choose the compact pair-of-pants core $P\subset B^\circ$ and
the zero-section whiskers used in \cref{prop:geometric-pi1}, and put
$J_P=f^{-1}(P)$.  Orient its boundary meridians clockwise.  Write
$\rho_0$ for the projection of the zero-section cusp meridian
$\widetilde a_0$.  The constructed candidate uses the tautological cusp
gluing, whose affine regluing class is $\mu=0$ (equivalently
$\ell_0=\gamma(\mu)=0$); hence the flat-marked cusp lift has no vertical
translation.  The base loops and their chosen lifts therefore satisfy
\[
 \rho_1\rho_2\rho_0=1,
 \qquad
 \widehat\rho_1\widehat\rho_2\widetilde a_0=1.
\tag{5.17g}
\]
For $L\in\operatorname{Pic}(X)$, write
$E=c_1(L|_F)\in\bigwedge^2V$.  Then $E\in(\bigwedge^2V)^G$, and there is a
crossed homomorphism
\[
 z:\pi_1(P)\longrightarrow V_{\mathbb R},\qquad
 z(gh)=z(g)+T_gz(h),
\tag{5.17h}
\]
whose class records
$c_1(L|_{J_P})-[\Omega_E]$.  Here $\Omega_E$ is the closed, translation-
invariant vertical two-form induced by $E$ in the flat marking.  If
$d_i=z(\rho_i)$ and a bar denotes the image in
$V_G\otimes\mathbb R$, then
\[
 \bar d_1+\bar d_2+\bar d_0=0.
\tag{5.17i}
\]
The three barred classes, and the number
$\langle d_0,\widehat\delta\rangle$, do not depend on the chosen crossed-
homomorphism representative.  More precisely, use the base-first convention
for the edge injection in (5.17j).  For any oriented loop $\rho\subset P$
and any $A_\rho$-invariant $\nu\in\Lambda$, its flat-marked sweep satisfies
\[
 \bigl\langle z(\rho),\nu\bigr\rangle
 =-\bigl\langle c_1(L|_{J_P})-[\Omega_E],
 a_*([S^1_\nu]\times[\widetilde\rho])\bigr\rangle.
\tag{5.17j'}
\]
Here $a$ is the fibrewise translation action and $\widetilde\rho$ is the
flat-marked lift of $\rho$.
\end{lemma}

\begin{proof}
The flat marking identifies the pullback of $J_P$ to the universal cover of
$P$ with
$\widetilde P\times(\Lambda_{\mathbb R}/\Lambda)$.  Since $E$ is the
restriction of a class on the total space, it is fixed by every monodromy
$T_g$.  The constant vertical form representing $E$ is therefore invariant
under the deck action and descends to the asserted form $\Omega_E$ on
$J_P$.

The pair of pants retracts onto a graph.  Hence the Serre spectral sequence
of the marked real torus bundle has only columns zero and one and gives the
short exact edge sequence
\[
0\longrightarrow H^1(P;V_{\mathbb R})
 \longrightarrow H^2(J_P;\mathbb R)
 \longrightarrow H^0(P;\textstyle\bigwedge^2V_{\mathbb R})
 \longrightarrow0.
\tag{5.17j}
\]
Both $c_1(L|_{J_P})$ and $[\Omega_E]$ map to $E$, so their difference is
the image of a unique class in $H^1(P;V_{\mathbb R})$.  The standard group-
cohomology model for a graph identifies this group with crossed
homomorphisms modulo principal crossed homomorphisms
\cite[Chapter~I, \S5]{Brown1982}; this gives (5.17h).

For completeness, the slant convention in (5.17j') is part of this
identification, not an additional geometric assertion.  Give an oriented
edge of a graph spine of $P$ a coordinate $s$ increasing with its based
loop and represent the corresponding $V_{\mathbb R}$-valued cellular
one-cochain by a base-first form $ds\wedge\beta_s$.  On the sweep of an
$A_\rho$-invariant fibre circle $S^1_\nu$, Fubini's formula gives
\[
\int_{S^1_s\times S^1_\nu}ds\wedge\beta_s
 =\int_{S^1_s}\langle\beta_s,\nu\rangle,ds
 =\langle z(\rho),\nu\rangle.
\]
The sweep in (5.17j') is written in the fibre-first order
$[S^1_\nu]\times[\widetilde\rho]$, which differs from this base-first
cross-product orientation by $(-1)^{1\cdot1}=-1$; this is the displayed
minus sign.
Changing the cellular representative by a principal cochain changes both
sides by the same boundary term, which vanishes because $A_\rho\nu=\nu$.
This proves (5.17j').

Applying $z$ to (5.17g) gives
\[
 d_1+T_1d_2+T_1T_2d_0=0.
\]
All $T_g$ act trivially on the global coinvariants, proving (5.17i).  A
change of representative has the form
$z(g)\mapsto z(g)+(T_g-I)a$.  Its image in $V_G\otimes\mathbb R$ is zero.
Moreover the displayed matrices in
\cref{prop:monodromy-certificate} fix $\widehat\delta$ under every dual
action $A_g$.  Therefore
\[
 \langle(T_g-I)a,\widehat\delta\rangle
 =\langle a,(A_g^{-1}-I)\widehat\delta\rangle=0,
\]
which proves the final assertion.
\end{proof}

\begin{lemma}[the finite-cover norm identity]
\label{lem:picard-finite-norm}
\Verified\ Retain the notation of
\cref{prop:candidate-finite-fillings,lem:picard-crossed-defect}.  For
$j=1,2$, let
\[
 N_{T_j}=I+T_j+\cdots+T_j^{m_j-1}.
\]
With the clockwise conventions of (5.17g), the pullback to the
$m_j$-fold $s_j$-cover gives
\[
 N_{T_j}d_j=\iota_{v_j}E
 \quad\text{in }V_{\mathbb R},
\qquad
 m_j\bar d_j=[\iota_{v_j}E]
 \quad\text{in }V_G\otimes\mathbb R.
\tag{5.17k}
\]
\end{lemma}

\begin{proof}
Restriction of the crossed homomorphism to the subgroup generated by
$\rho_j^{m_j}$ has value
\[
 z(\rho_j^{m_j})
 =\sum_{r=0}^{m_j-1}T_j^rz(\rho_j)=N_{T_j}d_j.
\tag{5.17l}
\]
This is the cohomological pullback of the annular defect to the
$s_j$-cover; the norm, rather than unqualified multiplication by $m_j$, is
essential before passage to global coinvariants.  A principal change
$d_j\mapsto d_j+(T_j-I)a$ does not alter this value, since
$N_{T_j}(T_j-I)a=(T_j^{m_j}-I)a=0$.

Let $Q_j:J'_j\to N_j$ be the finite quotient map.  The global flat marking
$\Xi$ over the full disc $\Delta_j$ makes the constant vertical form
$\Omega_E$ well defined on $J'_j$.  Radial retraction onto the central
torus makes
\[
 H^2(J'_j;\mathbb R)\longrightarrow H^2(A_0^{(j)};\mathbb R)
\]
an isomorphism.  On the central fibre the quotient is exactly $\pi_j$, and
naturality gives
\[
 (Q_j|_{A_0^{(j)}})^*c_1(L|_{S_j})
 =\pi_j^*c_1(L|_{S_j})=E.
\tag{5.17l'}
\]
The second equality is the actual radial/stalk identification in
\cref{lem:finite-filling-stalk-map}; on a nearby fibre the gluing map $h_j$
is a translation and therefore acts trivially on $H^2$ of that fibre.
Consequently the two classes
$c_1(Q_j^*(L|_{N_j}))$ and $[\Omega_E]$ have the same central restriction,
and the retraction isomorphism proves
\[
 c_1(Q_j^*(L|_{N_j}))-[\Omega_E]=0
 \quad\text{in }H^2(J'_j;\mathbb R).
\tag{5.17l''}
\]

On $\Delta_j^*$, the source identity
$h_j\widetilde g_j^{\log}h_j^{-1}=\widetilde g_j^0$ shows that $h_j$
descends from the logarithmic quotient $N_j^*$ to the linear/global
quotient $J^*$.  Thus
\[
 Q_j^*(L|_{N_j})|_{\Delta_j^*}
 \simeq h_j^*(L|_{J^*}).
\tag{5.17l'''}
\]
The same $h_j$ is translation in flat coordinates by $\alpha v_j$, where
\[
 \alpha=\frac{\log s_j}{2\pi i}.
\]
For any constant vertical two-form, direct substitution gives
\[
 h_j^*\Omega_E-\Omega_E
 =d\alpha\wedge\iota_{v_j}E.
\tag{5.17m}
\]
Along a clockwise loop in the $s_j$-plane one has
$\int d\alpha=-1$.  Thus the cross component contributed by (5.17m) is
$-\iota_{v_j}E$.  The remaining cross component is precisely (5.17l),
while the total class extends over the disc and has zero cross component.
Therefore $N_{T_j}d_j-\iota_{v_j}E=0$, which is the first identity in
(5.17k).  In $V_G$ every $T_j$ is the identity, so the norm becomes
multiplication by $m_j$ and the second identity follows.
\end{proof}

\begin{lemma}[the cusp boundary defect and Chern-class specialization]
\label{lem:picard-cusp-defect}
\Verified\ In the notation of
\cref{lem:picard-crossed-defect}, one has
\[
 \bar d_0=0\quad\text{in }V_G\otimes\mathbb R.
\tag{5.17n}
\]
Moreover, for every line bundle $L$ on the cusp neighbourhood $N_0$,
nearby-cycle specialization is compatible with first Chern classes:
\[
 \operatorname{sp}^2\bigl(c_1(L|_W)\bigr)=c_1(L|_F).
\tag{5.17o}
\]
\end{lemma}

\begin{proof}
Let $\widetilde a_0$ be the flat-marked clockwise zero-section meridian in
the lifted relation of (5.17g), and let
\[
 \xi_{\widehat\delta}
 =a_*([S^1_{\widehat\delta}]\times[\widetilde a_0])
\]
be its sweep under fibrewise translation by $\widehat\delta$.  The boundary
$\rho_0$ of $P$ was chosen to be its projection.  This is the zero-section
sweep: there is no hidden vertical drift in $\widetilde a_0$, by the
explicit overlap calculation in \cref{prop:geometric-pi1}.  Hence the
vertical form $\Omega_E$ vanishes on the swept torus, and the slant identity
(5.17j') gives
\[
 \langle d_0,\widehat\delta\rangle
 =-\langle c_1(L|_{J_P})-[\Omega_E],
          \xi_{\widehat\delta}\rangle
 =-\langle c_1(L|_{J_P}),\xi_{\widehat\delta}\rangle.
\tag{5.17p}
\]
Combine the exponential overlap coordinates $(x_1,x_2,t)$ of (2.8ii) with
the unimodular cusp chart
\[
 (u,x_1,x_2)=\left(\frac{t}{x_1x_2},x_1,x_2\right).
\]
The literal solid torus
\[
 \mathbb S_{\widehat\delta}
 =\{|u|=\rho,\ x_1=1,\ |x_2|\le1\}
\]
has boundary parametrization
\[
 (\theta,\phi)\longmapsto
 \bigl(\rho e^{i(\theta-\phi)},1,e^{i\phi}\bigr)
\tag{5.17p0}
\]
and hence $t=\rho e^{i\theta}$; its $\phi$-circle is translation by
$\widehat\delta$.  Direct evaluation in the proved cusp marking therefore
identifies $\partial\mathbb S_{\widehat\delta}$ with
$\xi_{\widehat\delta}$.  Reversing $\theta$ supplies the fixed clockwise
convention and only changes the common orientation sign.  No reduced
Mayer--Vietoris block or Leray comparison is used here.
Since $L$ extends across $N_0$, its Chern class evaluates to zero on that
boundary.  Thus the left side of (5.17p) is zero.  By
\cref{thm:coinvariants}, $V_G=\mathbb Z[\delta]$, and evaluation on the
globally invariant dual vector $\widehat\delta$ identifies
$V_G\otimes\mathbb R$ with $\mathbb R$.  This proves (5.17n).

For (5.17o), let $\ell=c_1(L|_{N_0})$, let
$\iota_0:W\hookrightarrow N_0$, and let
$\iota_t:F_t\hookrightarrow N_0$ be a sufficiently near smooth cusp fibre.
The neighbourhood realization of the nearby-cycle unit
$\mathbb Z_W\to\psi_{f_0}\mathbb Z$ recalled in the proof of
\cref{thm:p0-integral-specialization} gives the commutative square
\[
\begin{CD}
 H^2(N_0;\mathbb Z) @>{\iota_0^*}>> H^2(W;\mathbb Z)\\
 @V{\iota_t^*}VV @VV{\operatorname{sp}^2}V\\
 H^2(F_t;\mathbb Z) @= H^2(F_t;\mathbb Z).
\end{CD}
\tag{5.17p'}
\]
Applying this square to $\ell$ gives
$\operatorname{sp}^2(c_1(L|_W))=c_1(L|_{F_t})$.  Flat parallel transport
in the marked smooth family identifies the latter with $c_1(L|_F)=E$.
Equivalently, $c_1$ is the connecting morphism in the exponential sequence,
and restriction to the shrinking neighbourhoods and to $F_t$ commutes with
that morphism before passage to the direct limit.  This proves the claimed
naturality without using any global Leray differential.
\end{proof}

\begin{proposition}[direct boundary-Picard vanishing]
\label{prop:direct-boundary-picard}
\Verified\ For the constructed threefold and every analytic line bundle
$L\in\operatorname{Pic}(X)$,
\[
 c_1(L|_F)=0,
 \qquad
 \eta^*i^*L\simeq\mathcal O_{\mathrm{dP}_6}.
\tag{5.17q}
\]
This proof is independent of the global Leray transgressions
$d_2^{0,q}$ and of the packaged conclusion $H^2(X;\mathbb Z)=0$; it uses
only the actual finite-filling lattices, their logarithmic overlap maps, the
based boundary relation, the local cusp solid torus, and the local integral
specialization theorem.
\end{proposition}

\begin{proof}
The common invariant calculation in \cref{thm:coinvariants} first gives
\[
 E=c_1(L|_F)=kq.
\]
Because $L$ extends across the order-four filling, its fibre class belongs
to the actual specialization lattice
\[
 \pi_2^*H^2(S_2;\mathbb Z)
 =\{a\gamma\eta_2+bq:a\equiv b\pmod2\}
\]
of \cref{prop:finite-indices,lem:finite-filling-stalk-map}.  Taking
coefficient independence in the displayed invariant basis forces
$a=0$ and $b=k$; hence $k$ is even.  Write $E=2nq$.

Contraction with the two affine translation vectors is an exact integral
calculation:
\[
 \iota_{v_1}(2q)=4\eta_1,
 \qquad
 \iota_{v_2}(2q)=-6\eta_2.
\tag{5.17r}
\]
Indeed $v_1=(1,2,-4,0)^t$, $v_2=(-1,-3,3,0)^t$ in the basis dual to
$(\gamma,u,w,\delta)$, while
$q=uw+6\gamma\delta$.  In the global coinvariants,
\[
 [\eta_1]=3[\delta],\qquad [\eta_2]=2[\delta].
\]
Applying \cref{lem:picard-finite-norm} with $m_1=3$, $m_2=4$ therefore
gives
\[
 \bar d_1=4n[\delta],\qquad
 \bar d_2=-3n[\delta].
\tag{5.17s}
\]
Now combine the balance relation (5.17i) with the cusp vanishing (5.17n):
\[
 0=\bar d_1+\bar d_2+\bar d_0=n[\delta].
\]
The coinvariant group is torsion-free, so $n=0$ and hence $E=0$.

Finally (5.17o) and the injectivity of the local integral specialization in
\cref{thm:p0-integral-specialization} imply
$c_1(L|_W)=0$.  Pulling back to the normalization gives
$c_1(\eta^*i^*L)=0$.  The injectivity of
$c_1:\operatorname{Pic}(\mathrm{dP}_6)\to
H^2(\mathrm{dP}_6;\mathbb Z)$ proved in
\cref{prop:pic0-pullback} yields the second assertion of (5.17q).
\end{proof}

\subsection{An analytic completion of the all-Picard certificate}
\label{subsec:analytic-all-picard}

The local example in \cref{prop:local-cusp-picard-obstruction} shows that
the cusp neighbourhood alone cannot annihilate normalized central
restrictions.  We now close the stronger analytic route by retaining the
order-three and order-four fillings.  The proof has two independent parts.
First, logarithms of the actual \(\mathcal O^*\)-factors give the boundary
cocycle and its two finite-filling norm identities.  Second, a
candidate-specific support-function calculation reads every possible
central restriction directly from the mixed Appell--Humbert block on a
nearby torus.  No global Leray differential, nearby-cycle specialization,
or Chern-class evaluation on a cusp bounding chain occurs.

\begin{proposition}[analytic Appell--Humbert boundary cocycle]
\label{prop:analytic-ah-boundary}
\Verified\ Let \(P\subset B^\circ\) be the compact pair-of-pants core,
orient its boundary meridians clockwise, and retain the zero-section lifts
of \cref{prop:geometric-pi1}.  Thus
\[
 \rho_1\rho_2\rho_0=1,
 \qquad
 \widehat\rho_1\widehat\rho_2\widetilde a_0=1,
\tag{5.17-A1}
\]
and the second relation has no terminal vertical translation.  For
\(L\in\operatorname{Pic}(X)\), let
\[
 E=c_1(L|_F)\in\operatorname{Alt}^2(\Lambda,\mathbb Z)
   =\bigwedge^2V
\]
be its fibrewise Appell--Humbert form.  Put
\(V_{\mathbb R}=\operatorname{Hom}(\Lambda,\mathbb R)\), and use
\[
 A_g=(T_g^{-1})^{\mathsf t},
 \qquad
 \langle T_g\alpha,\lambda\rangle
 =\langle\alpha,A_g^{-1}\lambda\rangle.
\tag{5.17-A2}
\]
There is a crossed homomorphism, unique up to a principal crossed
homomorphism,
\[
 d:\pi_1(P)\longrightarrow V_{\mathbb R},
 \qquad
 d(gh)=d(g)+T_gd(h),
 \qquad
 d(g)\sim d(g)+(T_g-I)a.
\tag{5.17-A3}
\]
If \(d_i=d(\rho_i)\), then
\[
 d_1+T_1d_2+T_1T_2d_0=0,
 \qquad
 \bar d_1+\bar d_2+\bar d_0=0
 \quad\text{in }V_G\otimes\mathbb R.
\tag{5.17-A4}
\]
For the finite filling of order \(m_j\), \(j=1,2\),
\[
 N_{T_j}d_j=\iota_{v_j}E,
 \qquad
 m_j\bar d_j=[\iota_{v_j}E],
 \qquad
 N_{T_j}=I+T_j+\cdots+T_j^{m_j-1},
\tag{5.17-A5}
\]
where \((\iota_vE)(\lambda)=E(v,\lambda)\).  At the cusp,
\[
 \langle d_0,\widehat\delta\rangle=0,
 \qquad
 \bar d_0=0\quad\text{in }V_G\otimes\mathbb R.
\tag{5.17-A6}
\]
These identities are obtained from holomorphic transition factors and the
argument principle.
\end{proposition}

\begin{proof}
Choose an open Riemann-surface neighbourhood \(P^{\rm an}\Subset B^\circ\)
which deformation retracts onto \(P\), and pull the period family to its
universal cover.  Before quotienting by the fibre lattice, the resulting
space is
\[
 \widetilde{\mathcal J}
 =\widetilde P^{\rm an}\times\mathbb C^2.
\]
It is Stein and contractible, so the pullback of \(L\) is holomorphically
trivial.  In a nowhere-zero frame, write the factor belonging to a deck
transformation \(\alpha\) as
\[
 a_\alpha\in\mathcal O^*(\widetilde{\mathcal J}),
 \qquad
 a_{\alpha\beta}(x)=a_\alpha(\beta x)a_\beta(x).
\tag{5.17-A7}
\]
The space is simply connected, hence \(a_\alpha=e^{2\pi i h_\alpha}\) for
a holomorphic logarithm.  Therefore
\[
 c(\alpha,\beta)
 =h_\alpha(\beta x)+h_\beta(x)-h_{\alpha\beta}(x)\in\mathbb Z
\tag{5.17-A8}
\]
is independent of \(x\), and associativity makes \(c\) an integral group
two-cocycle.

Let \(\tau_\lambda\) denote fibre translation by
\(\Pi(z)\lambda\), \(\lambda\in\Lambda\), and let \(\widehat g\) be the
zero-section lift of a base deck transformation.  Period equivariance gives
\[
 \widehat g\,\tau_\lambda\,\widehat g^{-1}
 =\tau_{A_g\lambda}.
\tag{5.17-A9}
\]
On fibre translations, the commutator of (5.17-A8) is
\[
 E(\lambda,\mu)
 =c(\tau_\lambda,\tau_\mu)-c(\tau_\mu,\tau_\lambda).
\tag{5.17-A10}
\]
The four oriented sides of the \((\lambda,\mu)\)-parallelogram identify
this integer with the exponential connecting class of \(L|_F\); this is
the Appell--Humbert convention used in
(5.17av)--(5.17ax), and agrees with the standard factor-of-automorphy
calculation \cite[Chapter~1, \S2, pp.~5--8, especially
Theorem~2.1, p.~7]{BirkenhakeLange1999}%
\cite[equations~(3.3)--(3.10), pp.~3--7]{SchnellLecture3}.
Conjugating the commutator by \(\widehat g\) yields
\[
 E(A_g\lambda,A_g\mu)=E(\lambda,\mu).
\tag{5.17-A11}
\]

Push the logarithmic central extension defined by \(c\) out along
\(\mathbb Z\hookrightarrow\mathbb R\), and denote its central element of
weight \(r\) by \(\mathbf c(r)\).  Choose lifts \(U_\lambda\) of the fibre
translations satisfying the real Appell--Humbert normalization
\[
 U_\lambda U_\mu
 =\mathbf c\!\left(\frac12E(\lambda,\mu)\right)U_{\lambda+\mu}.
\tag{5.17-A12}
\]
To see that this normalization exists, start with arbitrary lifts.  The
difference between their real cocycle and \(E/2\) has zero commutator, so
it defines an abelian extension of the free abelian group \(\Lambda\) by
\(\mathbb R\).  Choose lifts of an integral basis and extend
multiplicatively; this splits that abelian extension and gives (5.17-A12).
If \(U'_\lambda\) is a second normalized system, then
\(U'_\lambda U_\lambda^{-1}\) is central, say
\(\mathbf c(\langle a,\lambda\rangle)\).  Comparing the two copies of
(5.17-A12) proves additivity in \(\lambda\), so
\(a\in V_{\mathbb R}\).  Thus every two normalized systems differ by exactly
the character change used below.

For \(g\in\pi_1(P)\), choose a lift \(G_g\) of \(\widehat g\) and define
\(d(g)\in V_{\mathbb R}\) by
\[
 G_gU_{A_g^{-1}\lambda}G_g^{-1}
 =\mathbf c\bigl(\langle d(g),\lambda\rangle\bigr)U_\lambda.
\tag{5.17-A13}
\]
The exponent is additive in \(\lambda\): conjugate (5.17-A12) and use
(5.17-A11).  The lifts of \(gh\) and \(G_gG_h\) differ centrally, so
\[
\begin{aligned}
 \langle d(gh),\lambda\rangle
 &=\langle d(g),\lambda\rangle
   +\langle d(h),A_g^{-1}\lambda\rangle\\
 &=\langle d(g)+T_gd(h),\lambda\rangle.
\end{aligned}
\]
This proves the crossed-homomorphism law.  Replacing
\(U_\lambda\) by
\(\mathbf c(\langle a,\lambda\rangle)U_\lambda\) changes (5.17-A13) by
\((T_g-I)a\).  Replacing \(G_g\) by a central multiple does not change
(5.17-A13), so this proves the stated independence and uniqueness.  Applying
the crossed law to
(5.17-A1) proves the first identity in (5.17-A4); passage to coinvariants
proves the second.

For the order-\(m_j\) filling, distinguish every covering level.  Put
\[
 \widetilde J_j=\Delta_j\times\mathbb C^2,\qquad
 r_j:\widetilde J_j\longrightarrow J'_j
\]
for the lattice quotient, whose deck translations are \(\tau_\lambda\).
On the vector cover define
\[
 Q_{j,\log}:=\widetilde g_j^{\,v_j/m_j},\qquad
 Q_{j,\mathrm{lin}}:=\widetilde g_j^{\,0}.
\]
They satisfy
\[
\begin{aligned}
 Q_{j,\log}\tau_\lambda Q_{j,\log}^{-1}
  &=\tau_{A_j\lambda},&
 (Q_{j,\log})^{m_j}&=\tau_{v_j},\\
 Q_{j,\mathrm{lin}}\tau_\lambda Q_{j,\mathrm{lin}}^{-1}
  &=\tau_{A_j\lambda},&
 (Q_{j,\mathrm{lin}})^{m_j}&=1.
\end{aligned}
\tag{5.17-A14}
\]
The positive sign is fixed by
\(s_j\circ g_j=e^{-2\pi i/m_j}s_j\) and (2.8f): after exactly
\(m_j\) clockwise iterations, the logarithmic vector-cover lift translates
by \(+v_j\).

Let \(q_{j,\log},q_{j,\mathrm{lin}}\) be the descended automorphisms of \(J'_j\),
and put \(J_j^{\prime *}=J'_j|_{\Delta_j^*}\).  Both descended automorphisms have
order dividing \(m_j\), because \(\tau_{v_j}\) is trivial after the lattice
quotient; each covers \(g_j\), whose order is exactly \(m_j\), so each
descended order is exactly \(m_j\).
The branch-independent torus translation
\[
 h_j:J_j^{\prime *}\longrightarrow J_j^{\prime *}
\]
constructed in (2.9d) satisfies, on this quotient level,
\[
 h_jq_{j,\log}h_j^{-1}=q_{j,\mathrm{lin}}.
\]
Define
\[
\begin{aligned}
 p_{j,\log}:J_j^{\prime *}&\longrightarrow
  N_j^*:=J_j^{\prime *}/\langle q_{j,\log}\rangle,\\
 p_{j,\mathrm{lin}}:J_j^{\prime *}&\longrightarrow
 J_j^*:=J_j^{\prime *}/\langle q_{j,\mathrm{lin}}\rangle .
\end{aligned}
\]
The conjugacy induces a biholomorphism
\[
 \bar h_j:N_j^*\longrightarrow J_j^*,\qquad
 p_{j,\mathrm{lin}}\circ h_j=\bar h_j\circ p_{j,\log}.
\]
The gluing of \(X\) identifies
\(L|_{N_j^*}\simeq\bar h_j^*(L|_{J_j^*})\); hence the exactly typed
line-bundle square gives
\[
 p_{j,\log}^*(L|_{N_j^*})
 \simeq h_j^*p_{j,\mathrm{lin}}^*(L|_{J_j^*}).
\tag{5.17-A14a}
\]

For the factor extension return to the filled vector cover.  Let
\[
 P_{j,\log}:J'_j\longrightarrow N_j,\qquad
 \pi_j^{\mathrm{vec}}:=P_{j,\log}\circ r_j:
 \widetilde J_j\longrightarrow N_j .
\]
Here \(P_{j,\log}\) is the full quotient and
\(P_{j,\log}|_{J_j^{\prime *}}=p_{j,\log}\).
The deck group is
\[
 \Gamma_j^{\log}
 =\langle\tau_\lambda,Q_{j,\log}:\lambda\in\Lambda\rangle.
\]
The pullback \((\pi_j^{\mathrm{vec}})^*(L|_{N_j})\) is trivial on the
contractible Stein manifold \(\widetilde J_j\); its factors define an
integral central extension, which we push out along
\(\mathbb Z\hookrightarrow\mathbb R\).  Over \(\Delta_j^*\), the
line-bundle square (5.17-A14a) identifies each fibre with the global factor
extension by a translation, so its commutator form is \(E\).  That form is
integer-valued and locally constant on the connected disc \(\Delta_j\);
therefore it is \(E\) on the filled cover as well.  Choose vertical lifts
\(U_\lambda^{(j)}\) normalized by
\[
 U_\lambda^{(j)}U_\mu^{(j)}
 =\mathbf c\!\left(\frac12E(\lambda,\mu)\right)
  U_{\lambda+\mu}^{(j)}.
\]
For \(\alpha\in\Gamma_j^{\log}\), define \(A_\alpha,T_\alpha\) by
\[
 \alpha\tau_\lambda\alpha^{-1}=\tau_{A_\alpha\lambda},\qquad
 T_\alpha=(A_\alpha^{-1})^{\mathsf t}.
\]
For a lift \(\widetilde G_\alpha\), define
\(\widetilde d_j(\alpha)\in V_{\mathbb R}\) through
\[
 \widetilde G_\alpha U_{A_\alpha^{-1}\lambda}^{(j)}
 \widetilde G_\alpha^{-1}
 =\mathbf c\!\left(\langle\widetilde d_j(\alpha),\lambda\rangle\right)
  U_\lambda^{(j)}.
\tag{5.17-A14c}
\]
Indeed, before naming \(\widetilde d_j(\alpha)\), the quotient of the two
sides of (5.17-A14c) is central because both cover the same translation.
Conjugating the normalized product and using
\[
 E(A_\alpha\lambda,A_\alpha\mu)=E(\lambda,\mu)
\]
shows that its exponent \(r_\alpha(\lambda)\) satisfies
\[
 r_\alpha(\lambda+\mu)=r_\alpha(\lambda)+r_\alpha(\mu).
\]
Thus it is the pairing with a unique element of \(V_{\mathbb R}\), as
asserted.  Since \(A_{\alpha\beta}=A_\alpha A_\beta\), choose the lift of
\(\alpha\beta\) to be \(\widetilde G_\alpha\widetilde G_\beta\); any other
choice differs centrally.  Then
\[
\begin{aligned}
 &\widetilde G_\alpha\widetilde G_\beta
 U_{A_\beta^{-1}A_\alpha^{-1}\lambda}^{(j)}
 \widetilde G_\beta^{-1}\widetilde G_\alpha^{-1}\\
 &\quad=
 \mathbf c\!\left(
  \langle\widetilde d_j(\beta),A_\alpha^{-1}\lambda\rangle
  +\langle\widetilde d_j(\alpha),\lambda\rangle\right)
 U_\lambda^{(j)}.
\end{aligned}
\]
The defining dual action
\(\langle T_\alpha v,\lambda\rangle
=\langle v,A_\alpha^{-1}\lambda\rangle\) therefore proves the crossed law
\[
 \widetilde d_j(\alpha\beta)
 =\widetilde d_j(\alpha)+T_\alpha\widetilde d_j(\beta).
\tag{5.17-A14d}
\]

Put
\[
 \widetilde J_j^*:=\Delta_j^*\times\mathbb C^2,
 \qquad
 r_j^*:=r_j|_{\widetilde J_j^*}:
 \widetilde J_j^*\longrightarrow J_j^{\prime *}
\]
and define
\[
\begin{aligned}
 \Phi_{j,\log}
  &:=p_{j,\log}\circ r_j^*,\\
 \Phi_{j,\mathrm{glob}}
  &:=p_{j,\mathrm{lin}}\circ h_j\circ r_j^*
   =\bar h_j\circ\Phi_{j,\log}.
\end{aligned}
\]
Both maps are invariant under \(\Gamma_j^{\log}\).  This is immediate for
every \(\tau_\lambda\), while
\[
\begin{aligned}
 \Phi_{j,\log}Q_{j,\log}
 &=p_{j,\log}q_{j,\log}r_j^*
  =\Phi_{j,\log},\\
 \Phi_{j,\mathrm{glob}}Q_{j,\log}
 &=p_{j,\mathrm{lin}}h_jq_{j,\log}r_j^*
  =p_{j,\mathrm{lin}}q_{j,\mathrm{lin}}h_jr_j^*
  =\Phi_{j,\mathrm{glob}}.
\end{aligned}
\]
Thus \(\widetilde J_j^*\) is the same \(\Gamma_j^{\log}\)-cover of the two
biholomorphic quotients.  Pulling the gluing isomorphism to this cover gives
\[
 \Theta_j:
 \Phi_{j,\log}^*(L|_{N_j^*})
 \xrightarrow{\ \sim\ }
 \Phi_{j,\mathrm{glob}}^*(L|_{J_j^*}).
\]
The filled factor extension restricts to this cover, and the global factor
extension pulls back to it through the fixed boundary whisker.  Because
\(\Theta_j\) intertwines descent, it induces an isomorphism of the resulting
pushed-out \(\Gamma_j^{\log}\)-factor extensions which sends
\(\mathbf c(r)\) to \(\mathbf c(r)\) for every \(r\in\mathbb R\).

Let \(\gamma_j^J:[0,1]\to J_j^*\) be the zero-section image of the
unwhiskered clockwise boundary meridian, and let \(\omega_j\) be the fixed
zero-section whisker from the global basepoint to \(\gamma_j^J(0)\).  Thus
the based total-space loop retained in (5.17-A1) is
\[
 \widehat\rho_j=\omega_j*\gamma_j^J*\omega_j^{-1}.
\]
Choose \(x_j\in J_j^{\prime *}\) such that
\(p_{j,\mathrm{lin}}(h_j(x_j))=\gamma_j^J(0)\), put
\(x_{j,\mathrm{lin}}=h_j(x_j)\), and let
\(\widetilde\gamma_{j,\mathrm{lin}}\) be the
\(p_{j,\mathrm{lin}}\)-lift of \(\gamma_j^J\) beginning at
\(x_{j,\mathrm{lin}}\).  For the chosen clockwise generator,
\[
 \widetilde\gamma_{j,\mathrm{lin}}(1)
 =q_{j,\mathrm{lin}}(x_{j,\mathrm{lin}}).
\]
Set
\[
 \gamma_j^N:=\bar h_j^{-1}\circ\gamma_j^J,
 \qquad
 \widetilde\gamma_{j,\log}
 :=h_j^{-1}\circ\widetilde\gamma_{j,\mathrm{lin}}.
\]
Then the quotient square gives the point-valued identities
\[
\begin{aligned}
 p_{j,\log}\circ\widetilde\gamma_{j,\log}
  &=\gamma_j^N,\\
 \widetilde\gamma_{j,\log}(0)&=x_j,\\
 \widetilde\gamma_{j,\log}(1)&=q_{j,\log}(x_j).
\end{aligned}
\]
Choose \(\widetilde x_j\in\widetilde J_j^*\) with
\(r_j^*(\widetilde x_j)=x_j\), and let
\(\widetilde\gamma_{j,\mathrm{vec}}\) be the \(r_j^*\)-lift of
\(\widetilde\gamma_{j,\log}\) beginning at \(\widetilde x_j\).  On
adjacent sheets choose compatible logarithm branches \(\ell,\ell'\) with
\[
 \ell'(w)=\ell(g_j^{-1}w)-\frac{2\pi i}{m_j}.
\]
The calculation in (2.9e) gives the point-valued endpoint formula
\[
 \widetilde\gamma_{j,\mathrm{vec}}(1)
 =Q_{j,\log}(\widetilde x_j),
\]
excluding both \(Q_{j,\log}^{-1}(\widetilde x_j)\) and every
\(\tau_\lambda Q_{j,\log}(\widetilde x_j)\) with \(\lambda\ne0\).
Fibre loops remain \(\tau_\lambda\), since \(h_j\) is a fibre translation.

The factor of \(L\) along the fixed lift of \(\omega_j\) identifies the
target factor over the local boundary segment with the globally based
factor.  Let \(G_{\rho_j}^{(j)}\) be the target lift of the local boundary
deck element whose whiskered image is the chosen global lift
\(G_{\rho_j}\).  Thus it covers \(\widehat\rho_j\) after whiskering and,
by (5.17-A13), satisfies
\[
 G_{\rho_j}^{(j)}U_{A_j^{-1}\lambda}(G_{\rho_j}^{(j)})^{-1}
 =\mathbf c(\langle d_j,\lambda\rangle)U_\lambda,
 \qquad d_j=d(\rho_j).
\]
Consequently, for some
\(a_j\in V_{\mathbb R}\) and \(b_j\in\mathbb R\),
\[
 \Theta_{j,*}(U_\lambda^{(j)})
 =\mathbf c(\langle a_j,\lambda\rangle)U_\lambda,
 \qquad
 \Theta_{j,*}(\widetilde G_{Q_{j,\log}})
 =\mathbf c(b_j)G_{\rho_j}^{(j)}.
\]
Here \(U_\lambda\) is the normalized global system in (5.17-A12), and
\(\mathbf c(b_j)\) disappears under conjugation.
Therefore
\[
\begin{aligned}
 \langle\widetilde d_j(Q_{j,\log}),\lambda\rangle
 +\langle a_j,\lambda\rangle
 &=\langle d_j,\lambda\rangle
   +\langle a_j,A_j^{-1}\lambda\rangle\\
 &=\langle d_j+T_ja_j,\lambda\rangle,
\end{aligned}
\]
so
\[
 \widetilde d_j(Q_{j,\log})=d_j+(T_j-I)a_j.
\tag{5.17-A14e}
\]

For a translation choose
\(\widetilde G_{\tau_v}=U_v^{(j)}\).  The filled normalization above gives
\[
 U_v^{(j)}U_\lambda^{(j)}(U_v^{(j)})^{-1}
 =\mathbf c(E(v,\lambda))U_\lambda^{(j)},
\]
and hence
\[
 \widetilde d_j(\tau_v)=\iota_vE,\qquad
 (\iota_vE)(\lambda)=E(v,\lambda).
\tag{5.17-A14f}
\]
Here
\[
 A_{Q_{j,\log}}=A_j,\qquad T_{Q_{j,\log}}=T_j.
\]
Since \(N_{T_j}(T_j-I)=T_j^{m_j}-I=0\), equations
(5.17-A14), (5.17-A14d), and (5.17-A14e) give
\[
\begin{aligned}
 N_{T_j}d_j
 &=N_{T_j}\widetilde d_j(Q_{j,\log})\\
 &=\widetilde d_j((Q_{j,\log})^{m_j})\\
 &=\widetilde d_j(\tau_{v_j})\\
 &=\iota_{v_j}E.
\end{aligned}
\tag{5.17-A15}
\]
Changing a line-bundle frame gives an isomorphic factor extension; changing
\(\widetilde G_\alpha\) by a central factor does not change conjugation; and
every second normalized vertical system differs by a character, changing
\(\widetilde d_j(\alpha)\) by \((T_\alpha-I)a\).  Finally,
\(\ell\mapsto\ell+2\pi in\) changes a vector-cover lift of \(h_j\) by
\[
 H_j\longmapsto H'_j=\tau_{nv_j}H_j.
\]
At the factor-extension level this means
\[
 \Theta'_{j,*}
 =\operatorname{Ad}(U_{nv_j})\circ\Theta_{j,*}.
\]
Consequently the comparison character changes by
\(n\iota_{v_j}E\), because
\[
 U_{nv_j}
 \bigl(\mathbf c(\langle a_j,\lambda\rangle)U_\lambda\bigr)
 U_{nv_j}^{-1}
 =
 \mathbf c(\langle a_j+n\iota_{v_j}E,\lambda\rangle)U_\lambda.
\]
Since \(A_jv_j=v_j\) and \(E\) is \(A_j\)-invariant,
\[
 \langle T_j(\iota_{v_j}E),\lambda\rangle
 =E(v_j,A_j^{-1}\lambda)
 =E(A_jv_j,\lambda)
 =E(v_j,\lambda),
\]
so this branch change leaves both (5.17-A14e) and (5.17-A15) unchanged.
The translation \(\tau_{nv_j}\) also commutes with
\(Q_{j,\mathrm{lin}}\) on the vector cover because \(A_jv_j=v_j\); on the
lattice quotient it is already the identity.  Hence it changes only the
central lift choice there.
The descended \(h_j,\bar h_j\), and the quotient square are branch
independent.  Thus (5.17-A15) is choice-independent and proves (5.17-A5),
including its sign.

It remains to obtain (5.17-A6) analytically.  Since
\(A_0\widehat\delta=\widehat\delta\), equation (5.17-A13) says that
\[
 [G_0,U_{\widehat\delta}]
 =\mathbf c\bigl(\langle d_0,\widehat\delta\rangle\bigr).
\tag{5.17-A16}
\]
This logarithmic commutator is the clutching winding with the base loop
first.  Thus, if \(b\) and \(f\) denote the oriented base and fibre loops,
respectively,
\[
 \operatorname{wind}_{b\wedge f}(L)
  =\langle d_0,\widehat\delta\rangle,
 \qquad
 \operatorname{wind}_{\xi_{\widehat\delta}}(L)
 =-\langle d_0,\widehat\delta\rangle.
\tag{5.17-A17}
\]
The second equality uses the fibre-first convention
\(\xi_{\widehat\delta}=[S^1_{\widehat\delta}]\times[\rho_0]\), hence the
single sign change.
Use the actual unimodular cusp chart
\[
 (u,x_1,x_2)=\left(\frac{t}{x_1x_2},x_1,x_2\right),
 \qquad t=ux_1x_2.
\tag{5.17-A18}
\]
For \(0<\rho\ll1\), the boundary parametrization
\[
 (\theta,\phi)\longmapsto
 \bigl(\rho e^{i(\theta-\phi)},1,e^{i\phi}\bigr),
 \qquad t=\rho e^{i\theta},
\tag{5.17-A19}
\]
has base coordinate \(\theta\) and fibre translation
\(\widehat\delta\) in the \(\phi\)-direction.  The latter fixed-base loop
does \emph{not} itself bound the \(x_2\)-disc.  To identify the correct
meridian, put \(\alpha=\theta-\phi\), let \(a_u\) be the positively
oriented \(u\)-longitude, and let \(m_{x_2}\) be the positively oriented
\(x_2\)-meridian at fixed \(u\).  In the basis
\((a_u,m_{x_2})\), the base and fibre loops are
\[
 b=a_u,
 \qquad
 f=m_{x_2}-a_u,
 \qquad
 b\wedge f=a_u\wedge m_{x_2}.
\tag{5.17-A19a}
\]
The change-of-basis determinant is \(+1\), so this preserves the
base-first orientation in (5.17-A17).  Now \(m_{x_2}\) does bound the
analytic disc
\[
 z\longmapsto(\rho e^{i\alpha_0},1,z),
 \qquad |z|\le1.
\tag{5.17-A19b}
\]

Choose a thin annulus \(A_u\) about \(|u|=\rho\), a disc \(D_{x_1}\)
about \(1\), and a disc \(D_{x_2}\) containing \(\{|x_2|\le1\}\), small
enough that
\[
 A_u\times D_{x_1}\times D_{x_2}
 \subset\{|ux_1x_2|<\epsilon\}.
\tag{5.17-A19c}
\]
This neighbourhood is Stein and retracts onto \(S^1\), so its second
integral cohomology vanishes.  Cartan's theorem~B and the exponential
sequence therefore give trivial Picard group
\cite[\S2, fundamental theorem, p.~81]{Cartan1957}%
\cite[Proposition~6.28 and equations~(6.35)--(6.36),
pp.~180--182]{Lee2024}.  Pull \(L|_{N_0}\) back to
the cusp cover and choose one holomorphic frame on this entire
neighbourhood.  Compute the clutching cocycle with this extending frame on
both sides of a cut in the \(u\)-longitude.  Its transition unit on
\(D_{x_2}\) is then identically \(F=1\), so by (5.17-A19a) its
base--fibre commutator winding is
\[
 \frac1{2\pi i}\int_{|x_2|=1}\frac{dF}{F}=0.
\tag{5.17-A20}
\]
This computes the same alternating logarithmic cocycle as (5.17-A16).
Indeed, changing from the global Appell--Humbert frame to the extending
chart frame changes \(c\) by a group coboundary.  Put
\[
 g=\widehat g_0,
 \qquad
 h=\tau_{\widehat\delta}
\]
in the underlying deck group.  Equation (5.17-A9) and
\(A_0\widehat\delta=\widehat\delta\) give \(gh=hg\).  If
\(c'=c+\delta\beta\) is the cocycle after a frame change, then
\[
 \bigl[c'(g,h)-c'(h,g)\bigr]
 -\bigl[c(g,h)-c(h,g)\bigr]
 =\beta(hg)-\beta(gh)=0.
\]
Thus the alternating cocycle is frame-invariant without assuming that its
central-extension lifts commute.  Equations (5.17-A17), (5.17-A19a), and
(5.17-A20) therefore give
\(\langle d_0,\widehat\delta\rangle=0\).  This is a calculation in
\(\mathcal O^*\), not an evaluation of a Chern class on a bounding chain.
Finally, the matrices \(A_1,A_2\) fix
\(\widehat\delta\), so the evaluation is unchanged by every principal
change in (5.17-A3).  Since
\(V_G=\mathbb Z[\delta]\) and
\(\delta(\widehat\delta)=1\), its vanishing is equivalent to
\(\bar d_0=0\).  This proves (5.17-A6).
\end{proof}

\begin{lemma}[the toric support function detects central restriction]
\label{lem:analytic-support-specialization}
\Verified\ Let \(M\in\operatorname{Pic}(N_0)\), let
\(\pi:Y_\epsilon\to N_0\) be the cusp covering, and let
\(\nu:E_0\to W\) be the normalization.  If the Appell--Humbert form of
\(M|_{F_t}\), \(t\ne0\), is \(E_M\), then the three independent degrees
of \(\nu^*(M|_W)\) determine an integral symmetric form
\(B:\mathbb Z^2\times\mathbb Z^2\to\mathbb Z\) satisfying
\[
 E_M\bigl((\lambda,0),(0,m)\bigr)
 =-B(B_0\lambda,m)
 \qquad(\lambda,m\in\mathbb Z^2).
\tag{5.17-A21}
\]
Consequently
\[
 E_M=0
 \quad\Longrightarrow\quad
 \nu^*(M|_W)\simeq\mathcal O_{E_0}.
\tag{5.17-A22}
\]
\end{lemma}

\begin{proof}
Put \(\mathcal M=\pi^*M\).  The identity
\(\pi\circ\Psi_\lambda=\pi\) gives \(\mathcal M\) its genuine
\(\Psi\)-equivariant descent structure.  In the notation of
\cref{prop:local-cusp-picard-obstruction}, order the six rays in the star
of the origin as
\[
 (r_1,r_2,r_3,r_4,r_5,r_6)
 =(e_1,e_2,e_2-e_1,-e_1,-e_2,e_1-e_2),
 \qquad C_i=E_0\cap E_{r_i}.
\tag{5.17-A23}
\]
The quotient identifies opposite conductor branches.  Therefore, for
\(A=\mathcal M|_{E_0}=\nu^*(M|_W)\),
\[
 d_i=\deg_{C_i}A,
 \qquad
 (\deg_{C_1}A,\ldots,\deg_{C_6}A)
 =(d_1,d_2,d_3,d_1,d_2,d_3).
\tag{5.17-A24}
\]

Define the integer-valued quadratic function
\[
 a(x,y)
 =(d_2+d_3)\binom{x}{2}+d_3xy+(d_1+d_3)\binom{y}{2}.
\tag{5.17-A25}
\]
Its polarization is
\[
 B(v,w)=a(v+w)-a(v)-a(w),
 \qquad
 [B]=
 \begin{pmatrix}
  d_2+d_3&d_3\\ d_3&d_1+d_3
 \end{pmatrix}.
\tag{5.17-A26}
\]
The locally finite divisor
\[
 D_a=\sum_{v\in\mathbb Z^2}a(v)E_v
\tag{5.17-A27}
\]
is Cartier: on every maximal chart, the three height-one ray generators
form an integral basis, so the three assigned coefficients are the values
of one integral character.  In the standard toric convention the associated
support function has ray values \(-a(v)\); the divisor coefficients here are
the values \(a(v)\), consistent with
\(\operatorname{div}(\chi^m)=\sum_\rho\langle m,u_\rho\rangle D_\rho\)
\cite[\S3.3, pp.~60--63; \S3.4, pp.~63--72]{Fulton1993}%
\cite[Propositions~4.1.1--4.1.2, p.~171; Theorems~4.2.8,
p.~179 and~4.2.12, p.~182]{CoxLittleSchenck2011}.
This is the same chartwise Cartier argument as
in (5.17ab), and it applies to an arbitrary integer-valued \(a\).

Direct substitution gives
\[
 (a(r_1),\ldots,a(r_6))
 =(0,0,d_2,d_2+d_3,d_1+d_3,d_1).
\tag{5.17-A28}
\]
Using \(C_i^2=-1\) and \(C_iC_{i\pm1}=1\), the degree vector of
\(\mathcal O(D_a)|_{E_0}\) is exactly the vector in (5.17-A24).  More
generally, restrict the principal identity
\(\operatorname{div}(t)=\sum_wE_w\) to \(E_v\).  Since \(E_v\) meets
precisely \(E_{v+r_i}\), it gives
\[
 \mathcal O(D_a)|_{E_v}
 \simeq
 \mathcal O_{E_v}\!\left(
   \sum_{i=1}^6(a(v+r_i)-a(v))C_{v,i}
 \right).
\tag{5.17-A29}
\]
The coefficient difference is
\[
 a(v+r_i)-a(v)=a(r_i)+B(v,r_i).
\tag{5.17-A30}
\]
The second summand in (5.17-A30) is the divisor of the toric character
with exponent \(B(v,-)\).  Thus every restriction in (5.17-A29) has the
same six degrees as \(\mathcal O(D_a)|_{E_0}\).

Because \(B_0\bar\Lambda=\mathbb Z^2\), the deck action is transitive on
the components \(E_v\).  The genuine descent structure on \(\mathcal M\)
therefore transports (5.17-A24), with its ordered six boundary directions,
to every \(E_v\).  It follows that
\[
 \mathcal R:=\mathcal M\otimes\mathcal O(-D_a)
 \quad\text{satisfies}\quad
 \mathcal R|_{E_v}\simeq\mathcal O_{E_v}
 \quad(v\in\mathbb Z^2).
\tag{5.17-A31}
\]
Indeed, the six toric curves detect \(\operatorname{Pic}(\mathrm{dP}_6)\):
in the blow-up basis \(H,E_1,E_2,E_3\), vanishing of the degrees on the
three exceptional curves kills the three \(E_i\)-coefficients, and
vanishing on any strict transform \(H-E_i-E_j\) kills the remaining
\(H\)-coefficient.  Since
\(H^1(\mathrm{dP}_6,\mathcal O)=0\), equality of those degrees implies
isomorphism, not merely numerical equivalence.

For \(b_\lambda=B_0\lambda\), the torus-translation part of
\(\Psi_\lambda\) preserves each toric prime divisor, while its fan part
translates the index by \(b_\lambda\).  Hence
\[
 \Psi_\lambda^*D_a-D_a
 =\sum_v(a(v+b_\lambda)-a(v))E_v
 =\operatorname{div}(\chi^{m_\lambda}),
\tag{5.17-A32}
\]
where
\[
 m_\lambda=(\ell_\lambda,k_\lambda),
 \qquad
 \ell_\lambda=B(b_\lambda,-),
 \qquad
 k_\lambda=a(b_\lambda).
\tag{5.17-A33}
\]
Thus multiplication by \(\chi^{m_\lambda}\) gives an individual
isomorphism
\(\Psi_\lambda^*\mathcal O(D_a)\simeq\mathcal O(D_a)\).
No cocycle assertion for these individual isomorphisms is needed: we use
them only to separate the horizontal winding of the genuine descent factor
of \(\mathcal M\).

For a ray \(\mathbb R_{\ge0}(v,1)\), its toric open set in
\(Y_\epsilon\) is
\[
 U_v\simeq\Delta_\epsilon\times(\mathbb C^*)^2,
 \qquad E_v^\circ=\{0\}\times(\mathbb C^*)^2.
\tag{5.17-A34}
\]
Indeed \((e_1,0),(e_2,0),(v,1)\) is an integral basis and \(t\) is the
character dual to the last vector; compare
\cite[\S2.1, pp.~28--29]{Fulton1993} and
\cite[Example~1.2.21, p.~31]{CoxLittleSchenck2011}.
The set \(U_v\) is Stein and retracts
onto \(E_v^\circ\).  By (5.17-A31),
\(c_1(\mathcal R|_{U_v})=0\); the exponential sequence and Cartan's
theorem~B therefore trivialize \(\mathcal R|_{U_v}\).  Choose a global
nowhere-zero frame \(s_v\) of \(\mathcal R|_{E_v}\).  A frame on \(U_v\) can be
multiplied by a unit pulled back from \((\mathbb C^*)^2\) so that its
restriction to \(E_v^\circ\) agrees with \(s_v|_{E_v^\circ}\); denote the
resulting frame by \(e_v\).

We also need to compare these frames on the generic torus.  The direction
of the toric affine-open inclusions is essential here.  If \(v,w\) are
adjacent vertices, choose a triangle \(\sigma=\{v,w,u\}\) containing them,
write \(z_v,z_w,z_u\) for the coordinates dual to its three height-one ray
generators, and put
\[
 S_\sigma:=U_\sigma\cap Y_\epsilon
 =\{(z_v,z_w,z_u)\in\mathbb C^3:
       |z_vz_wz_u|<\epsilon\}.
\tag{5.17-A34a}
\]
Because \(\rho_v<\sigma\), localization of the maximal affine chart gives
\begin{equation}
 U_{\rho_v}=D(z_wz_u)\subset U_\sigma,\qquad
 U_{\rho_w}=D(z_vz_u)\subset U_\sigma,\qquad
 U_{\rho_v}\cap U_{\rho_w}=D(z_vz_wz_u)=U_0.
\label{eq:analytic-toric-open-directions}
\end{equation}
Thus \(U_\sigma\) is not contained in \(U_{\rho_v}\cap U_{\rho_w}\), and
the coordinate discs in \(S_\sigma\) cannot be applied directly to the
ratio \(e_w/e_v\).

The space \(S_\sigma\) is nevertheless the required comparison chart.  It
is the closed graph
\[
 \{(z_v,z_w,z_u,t)\in\mathbb C^3\times\Delta_\epsilon:
   t=z_vz_wz_u\}
\]
in the Stein space \(\mathbb C^3\times\Delta_\epsilon\), and hence is
Stein.  Radial scaling of the three \(z\)-coordinates contracts it to the
origin.  The exponential sequence and Cartan's theorem~B therefore give
\(\operatorname{Pic}(S_\sigma)=0\).  Choose a frame \(e_\sigma\) of
\(\mathcal R|_{S_\sigma}\).

Let \(s_v\) be the global frame of
\(\mathcal R|_{E_v}\simeq\mathcal O_{E_v}\) used to normalize \(e_v\).
On \(U_{\rho_v}\cap Y_\epsilon\) put
\[
 A_{\sigma,v}=e_v/e_\sigma.
\]
At \(t=0\), this is the restriction to \((\mathbb C^*)^2\) of the
holomorphic unit
\[
 G_{\sigma,v}
 =\frac{s_v}{e_\sigma|_{E_v\cap S_\sigma}}
 \in\mathcal O^*(E_v\cap S_\sigma),
 \qquad E_v\cap S_\sigma=\{z_v=0\}\simeq\mathbb C^2.
\]
Consequently its winding on every horizontal cocharacter is zero at
\(t=0\).  Under
\(U_{\rho_v}\cap Y_\epsilon\simeq
\Delta_\epsilon\times(\mathbb C^*)^2\), that winding is an integer locally
constant in \(t\), and is therefore zero on every nearby nonzero fibre.
The same argument applies to \(A_{\sigma,w}\).  On their actual common
domain \(U_0\cap Y_\epsilon\),
\[
 e_w/e_v=A_{\sigma,w}/A_{\sigma,v},
\]
so this ratio has zero horizontal winding.  The one-skeleton of the
\(A_2\)-triangulation is connected, and multiplication along an edge path
now gives
\[
 \operatorname{wind}_m(e_w/e_v)=0
 \qquad(v,w,m\in\mathbb Z^2)
\tag{5.17-A34b}
\]
on every nonzero fibre.

Let
\(\theta^{\mathcal M}_\lambda:\Psi_\lambda^*\mathcal M\to\mathcal M\)
be the genuine descent isomorphism, let
\(\mu^+_\lambda:\Psi_\lambda^*\mathcal O(D_a)\to\mathcal O(D_a)\)
be multiplication by \(\chi^{m_\lambda}\) from (5.17-A32), and define the
inverse-dual isomorphism
\[
 \mu^-_\lambda:=((\mu^+_\lambda)^{-1})^\vee:
 \Psi_\lambda^*\mathcal O(-D_a)\longrightarrow\mathcal O(-D_a),
 \qquad s\longmapsto\chi^{-m_\lambda}s.
\]
The ordinary dual of \(\mu^+_\lambda\) has the opposite arrow; the
inverse-dual above is the map required by the pullback tensor product.  Then
\[
 \theta^{\mathcal R}_\lambda
 :=\theta^{\mathcal M}_\lambda\otimes\mu^-_\lambda:
 \Psi_\lambda^*\mathcal R\longrightarrow\mathcal R.
\]
Since \(\Psi_\lambda(U_v)=U_{v+b_\lambda}\), define
\(r_{\lambda,v}\in\mathcal O^*(U_v)\) by the typed identity
\begin{equation}
 \theta^{\mathcal R}_\lambda
   (\Psi_\lambda^*e_{v+b_\lambda})
 =r_{\lambda,v}e_v.
\label{eq:analytic-residual-descent-factor}
\end{equation}
At \(t=0\) it extends to an isomorphism between two trivial line bundles
on the compact surface \(E_v\).  Its factor in the chosen global frames is
therefore a nonzero constant.  The winding of
\(r_{\lambda,v}(t,-)\) around every horizontal one-parameter subgroup is
an integer locally constant in \(t\); at \(t=0\) it is zero, hence it is
zero for every nearby \(t\ne0\).

To express this isomorphism in one frame on the generic torus, put
\(q_{\lambda,v}=e_{v+b_\lambda}/e_v\).  Equation (5.17-A34b) gives zero
horizontal winding for \(q_{\lambda,v}\); because
\(\Psi_\lambda|_{Y_t}\) is multiplication by a fixed torus point, it
induces the identity on horizontal \(H_1\), and
\(\Psi_\lambda^*q_{\lambda,v}\) also has zero winding.  Relative to the
single frame \(e_v\), the residual multiplier is therefore
\[
 h_{\lambda,v}
 =\frac{r_{\lambda,v}}{\Psi_\lambda^*q_{\lambda,v}},
 \qquad
 \operatorname{wind}_m(h_{\lambda,v})=0.
\tag{5.17-A34c}
\]

On \(U_0\) let \(\mathbf1_a\) be the canonical toric frame of
\(\mathcal O(D_a)\), and put \(F_v=e_v\otimes\mathbf1_a\).  Equations
\cref{eq:analytic-residual-descent-factor} and (5.17-A32) give
\[
 \theta^{\mathcal M}_\lambda(\Psi_\lambda^*F_v)
 =h_{\lambda,v}\chi^{m_\lambda}F_v.
\]
On a nonzero fibre, \(\chi^{m_\lambda}=x^{\ell_\lambda}t^{k_\lambda}\).
With the section-descent convention (5.17am), the factor of automorphy is
therefore
\[
 \alpha_\lambda=(h_{\lambda,v}\chi^{m_\lambda})^{-1}.
\]
The residual factor has zero horizontal winding by (5.17-A34c), and hence
\begin{equation}
 \operatorname{wind}_m(\alpha_\lambda)
 =-\ell_\lambda(m)=-B(B_0\lambda,m).
\label{eq:analytic-factor-horizontal-winding}
\end{equation}

This winding is exactly the mixed Chern-class entry.  Pull \(F_v\) back to
the universal cover
\[
 p_t:\mathbb C^2\longrightarrow Y_t,
 \qquad \zeta\longmapsto(e^{2\pi i\zeta},t).
\]
The period indexed by \((\lambda,n)\in\mathbb Z^2\oplus\mathbb Z^2\)
acts by \(\zeta\mapsto\zeta+Z(s)\lambda+n\).  The horizontal factor is one,
because \(F_v\) is already defined on \(Y_t\).  Choose a logarithm
\(g_\lambda\in\mathcal O(\mathbb C^2)\) with
\[
 p_t^*\alpha_\lambda=e^{2\pi i g_\lambda}.
\]
For every horizontal period \(n\),
\[
 g_\lambda(\zeta+n)-g_\lambda(\zeta)
 =\operatorname{wind}_n(\alpha_\lambda).
\]
In the exponential connecting cocycle
\[
 c(l,l')=g_l(\zeta+l')+g_{l'}(\zeta)-g_{l+l'}(\zeta),
\]
take the horizontal logarithm to be zero and the logarithm for the combined
period \((\lambda,n)\) to be \(g_\lambda\).  Then
\[
 c((\lambda,0),(0,n))
 =\operatorname{wind}_n(\alpha_\lambda),
 \qquad
 c((0,n),(\lambda,0))=0.
\]
Therefore the alternating exponential connecting integer between
the vertical deck
translation indexed by \(\lambda\) and the horizontal loop indexed by
\(m\) is
\[
 E_M\bigl((\lambda,0),(0,m)\bigr)
 =-\ell_\lambda(m)
 =-B(B_0\lambda,m),
\]
which proves (5.17-A21).  This sign also recovers the explicit check in
\cref{prop:local-cusp-picard-obstruction}: for
\(a(v)=v^{\mathsf t}Hv\), one has \(B=2H\), so (5.17-A21) is the
cross-block \(-2B_0^{\mathsf t}H=-2P\), namely \(-2q\).

If \(E_M=0\), equation (5.17-A21) and the unimodularity of \(B_0\) force
\(B=0\).  The matrix in (5.17-A26) then gives
\(d_3=0\), \(d_2=0\), and \(d_1=0\).  Thus all six degrees of \(A\)
vanish, and the Picard detection argument following (5.17-A31) gives
\(A\simeq\mathcal O_{E_0}\).  This is (5.17-A22).
\end{proof}

\begin{theorem}[purely analytic vanishing of the global restriction]
\label{thm:analytic-all-picard}
\Verified\ For the constructed threefold,
\[
 r=\eta^*i^*:\operatorname{Pic}(X)\longrightarrow
 \operatorname{Pic}(\mathrm{dP}_6)
\]
is the zero homomorphism.  This proof uses only holomorphic
\(\mathcal O^*\)-factors, the two finite logarithmic fillings, the cusp
argument principle, and the \(\Psi\)-equivariant toric support function.
It uses neither a global integral Leray transgression nor local integral
nearby-cycle specialization.
\end{theorem}

\begin{proof}
Let \(L\in\operatorname{Pic}(X)\).  By
\cref{thm:coinvariants}, its invariant fibre class is
\[
 E=c_1(L|_F)=kq,
 \qquad k\in\mathbb Z,
 \qquad q=uw+6\gamma\delta.
\tag{5.17-A35}
\]
The exact contractions, now used without an integral specialization
lattice, are
\[
 \iota_{v_1}q=2\eta_1,
 \qquad
 \iota_{v_2}q=-3\eta_2,
 \qquad
 [\eta_1]=3[\delta],
 \qquad
 [\eta_2]=2[\delta].
\tag{5.17-A36}
\]
Apply (5.17-A5) with \((m_1,m_2)=(3,4)\).  In the real global
coinvariants,
\[
 \bar d_1=2k[\delta],
 \qquad
 \bar d_2=-\frac32k[\delta].
\tag{5.17-A37}
\]
Equations (5.17-A4) and (5.17-A6) therefore give
\[
 0=\bar d_1+\bar d_2+\bar d_0
   =\frac{k}{2}[\delta].
\tag{5.17-A38}
\]
Since \(V_G\otimes\mathbb R=\mathbb R[\delta]\), one has \(k=0\), so
\(E=0\).  Apply \cref{lem:analytic-support-specialization} to
\(M=L|_{N_0}\).  It gives
\[
 \eta^*i^*L=\nu^*(L|_W)\simeq\mathcal O_{E_0}.
\]
This holds for every \(L\), proving \(r=0\).
\end{proof}


\begin{corollary}[candidate nonvanishing for \texorpdfstring{$\operatorname{Pic}^0$}{Pic0} without topology]
\label{cor:candidate-pic0-route}
\Verified\ For the constructed threefold and every
$L\in\operatorname{Pic}^0(X)$,
\[
 H^0(W,\Omega_X^1|_W\otimes A_L)\ne0,
 \qquad
 H^2(W,(T_X\otimes L)|_W)\ne0,
\]
and
\[
 R^2f_*(T_X\otimes L)\ne0.
\tag{5.17f}
\]
This proof uses the analytic construction of the cusp fibre but no global
homology, Mayer--Vietoris, or Leray calculation.
\end{corollary}

\begin{proof}
The cusp theorem \cref{thm:candidate-cusp-geometry} supplies the full compact
reduced normal-crossings fibre, its $\mathrm{dP}_6$ normalization, and the
trivialization of $\omega_X$ near $W$.  By
\cref{prop:pic0-pullback}, $\eta^*A_L$ is trivial, so
\cref{prop:conductor} gives the first nonvanishing.  Apply
\cref{prop:cartier-duality-absolute} to $E=T_X\otimes L$; its dual space is
exactly the preceding nonzero section space, giving the fibrewise $H^2$
nonvanishing.  The fibre has complex dimension two, so its coherent
$H^3$ vanishes.  The global construction makes $X$ compact and $f$
nonconstant; \cref{lem:curve-flatness} therefore gives properness and
$f$-flatness of $T_X\otimes L$.  The base-change isomorphism of
\cref{prop:analytic-base-change} now identifies the fibre of the $R^2$
stalk with the nonzero $H^2$ group.  Hence the stalk, and therefore the
direct-image sheaf, is nonzero.
\end{proof}

\begin{remark}[reconciliation of the three restriction routes]
The Pic$^0$ argument proves (5.17f) without global homology,
Mayer--Vietoris, or Leray.  The all-Pic statement
\cref{prop:leray-edge-restriction} avoids citing the packaged theorem
$H^2(X;\mathbb Z)=0$, but it reuses the same verified global integral
Leray and specialization machinery from which that theorem is deduced.
By contrast, \cref{prop:direct-boundary-picard} supplies a second all-Pic
proof without any global Leray differential: it recovers the primitive
$4-3=1$ obstruction directly from the finite logarithmic gluing and the
cusp boundary.  It still uses the local integral cusp-specialization
theorem.  The third proof in \cref{thm:analytic-all-picard} replaces both
that specialization and the cusp Chern-class bounding-chain evaluation by
logarithmic $\mathcal O^*$ factors, the argument principle, and the exact
$\Psi$-equivariant support function.  Thus the purely analytic
strengthening is now complete; the three proofs have genuinely different
dependency sets.
\end{remark}

\subsection{Abstract conductor theorem for the candidate fibre}

\begin{theorem}[the candidate conductor mechanism, with all hypotheses]
\label{thm:candidate-conductor}
\Verified\ Assume:
\begin{enumerate}[label=(H\arabic*)]
\item $X$ is a connected compact smooth complex threefold and
$f:X\to B$ is a nonconstant holomorphic map to a smooth compact curve;
\item $W=f^{-1}(p_0)$ is the full scheme-theoretic fibre, is reduced, and
has the normal-crossings local models (5.3);
\item the normalization $\eta:\widetilde W\to W$ is the degree-six del
Pezzo surface $\mathrm{dP}_6$;
\item $H^2(X;\mathbb Z)=0$.
\end{enumerate}
Then, for every $L\in\operatorname{Pic}(X)$,
\[
H^2(W,(T_X\otimes L)|_W)\ne0
\quad\text{and}\quad
R^2f_*(T_X\otimes L)\ne0.
\tag{5.18}
\]
\end{theorem}

\begin{proof}
Fix $L$ and put
\[
\mathcal F_L=T_X\otimes L,\qquad
A_L=(L^\vee\otimes\omega_X)|_W.
\]
Hypothesis (H4) gives $c_1(L^\vee\otimes\omega_X)=0$.  Pullback therefore
gives $c_1(\eta^*A_L)=0$.  The surface $\mathrm{dP}_6$ is the blow-up of
$\mathbb P^2$ at three noncollinear points; the standard blow-up calculation
gives $H^1(\mathrm{dP}_6,\mathcal O)=0$
\cite[Chapter~V]{Hartshorne1977}.  In the exponential sequence, the kernel
of
\[
c_1:\operatorname{Pic}(\mathrm{dP}_6)\longrightarrow
H^2(\mathrm{dP}_6;\mathbb Z)
\]
is the image of $H^1(\mathrm{dP}_6,\mathcal O)$ and is therefore zero.
Consequently
\[
\eta^*A_L\simeq\mathcal O_{\widetilde W}.
\tag{5.19}
\]
\Cref{prop:conductor} now gives
\[
H^0(W,\Omega^1_X|_W\otimes A_L)\ne0.
\tag{5.20}
\]

Because $W=f^*(p_0)$, one has
$\mathcal O_X(W)=f^*\mathcal O_B(p_0)$, whose restriction to the fibre is
trivial.  Adjunction therefore gives $\omega_W\simeq\omega_X|_W$.
Apply \cref{thm:cartier-duality} to $E=\mathcal F_L$:
\[
\begin{split}
H^2(W,\mathcal F_L|_W)^\vee
&\simeq H^0(W,(\mathcal F_L|_W)^\vee\otimes\omega_W)\\
&=H^0(W,\Omega^1_X|_W\otimes A_L).
\end{split}
\]
The right side is nonzero by (5.20), proving the first assertion of (5.18).
The same duality theorem gives
$H^3(W,\mathcal F_L|_W)=0$.

Since $X$ is compact, $f$ is proper.  By
\cref{lem:curve-flatness}, $f$ is flat and the vector bundle
$\mathcal F_L$ is $f$-flat.  Thus every hypothesis of
\cref{prop:analytic-base-change} is satisfied at $p_0$, and
\[
(R^2f_*\mathcal F_L)_{p_0}\otimes\mathbb C(p_0)
\simeq H^2(W,\mathcal F_L|_W)\ne0.
\]
The stalk, and hence the coherent direct-image sheaf, is nonzero.  This is
the second assertion of (5.18).
\end{proof}

\begin{corollary}[application to the constructed threefold]
\label{cor:candidate-conductor-actual}
\Verified\ For the global construction of
\cref{thm:candidate-global-geometry} and every
$L\in\operatorname{Pic}(X)$,
\[
H^2(W,(T_X\otimes L)|_W)\ne0,
\qquad R^2f_*(T_X\otimes L)\ne0.
\]
\end{corollary}

\begin{proof}
The cusp construction gives the compact full Cartier fibre, its
$\mathrm{dP}_6$ normalization, and the local trivialization of $\omega_X$.
By the local canonical trivialization and
\cref{prop:direct-boundary-picard}, $\eta^*A_L$ is trivial for every $L$.
The conductor proposition then gives
$H^0(W,\Omega_X^1|_W\otimes A_L)\ne0$; absolute Cartier-fibre duality turns
this into the displayed $H^2$ nonvanishing, and the uniform top-degree
base-change isomorphism turns it into a nonzero $R^2$ stalk.  This is exactly
the chain proved in \cref{cor:candidate-pic0-route}, now with the all-Pic
restriction supplied directly by \cref{prop:direct-boundary-picard}.  The
global Leray-edge proof in \cref{prop:leray-edge-restriction} is an
independent certificate of the same restriction statement.
\end{proof}

\subsection{A second printed issue: abelianization of split extensions}

\begin{lemma}[abelianization of a semidirect product]
\label{lem:semidirect-abelianization}
\Verified\ If
$1\to K\to G\to Q\to1$ is split, then
\[
G^{\mathrm{ab}}\simeq
Q^{\mathrm{ab}}\oplus(K^{\mathrm{ab}})_Q,
\tag{5.21}
\]
where the subscript denotes coinvariants for the conjugation action of $Q$.
\end{lemma}

\begin{proof}
Choose a splitting and write $G=K\rtimes Q$.  Let $N\triangleleft K$ be
generated by $[K,K]$ and by all elements
$qkq^{-1}k^{-1}$ with $q\in Q$, $k\in K$.  It is normal in $G$:
$[K,K]$ makes it stable under conjugation by $K$, while conjugation by
$Q$ permutes the displayed cross-commutators.  By construction
$K/N=(K^{\mathrm{ab}})_Q$, and $Q$ acts trivially on this quotient.  Hence
\[
G/N\simeq(K^{\mathrm{ab}})_Q\times Q.
\]
Since $N\subseteq[G,G]$, abelianizing the last display gives (5.21).
\end{proof}

The sentence in \cite[Lemma~3.2, pp.~84--85]{CDP1998} that passes from a
split fundamental-group extension to a direct sum with $H_1(F;\mathbb Z)$
is therefore valid only when monodromy acts trivially after abelianization.
A minimal counterexample retaining a free quotient group is
\[
G=\mathbb Z\rtimes_{-1}\mathbb Z.
\]
The extension splits, but
$G^{\mathrm{ab}}\simeq\mathbb Z\oplus\mathbb Z/2$, whereas the printed
direct-product formula gives $\mathbb Z^2$.
For the candidate's printed monodromy, \cref{thm:coinvariants} gives
\[
H^1(F;\mathbb Q)^{\pi_1(B^\circ)}=\mathbb Q\gamma,
\]
so the rational coinvariant rank is one, not
$b_1(F)=4$.

If $C^\circ$ is a sphere punctured at $s$ points, the correct rank formula is
\[
b_1(X^\circ)=s-1+
\dim_{\mathbb Q}H_1(F;\mathbb Q)_{\pi_1(C^\circ)}.
\tag{5.22}
\]
Thus the proof of CDP98 Theorem~3.1 reaches only the coinvariant rank unless
an additional theorem proves trivial monodromy.  CDP20 Lemma~4.2 repeats the
direct-sum formula, and CDP20 Corollary~5.2 imports CDP98 Theorem~3.1.
Repairing the local-freeness issue discussed there does not repair this
distinct monodromy error.  The direct-product step in the torus branch of the
CDP20 proof is therefore refuted, independently of the singular-fibre defect
above.

\subsection{The corrected theorem that survives both defects}

The original hypotheses do not force the CDP conclusion: the constructed
threefold above is a counterexample.  There is nevertheless a sharp useful
replacement.  It isolates the exact extra global vanishing needed by the
intact part of the argument, and the following fibrewise criterion replaces
the invalid p.~692 reduction.

\begin{lemma}[generic endpoint vanishing]\label{lem:cdp-generic-endpoints}
\Verified\ Let \(Y\) be a connected compact complex manifold, let \(D\) be
a nonzero effective Cartier divisor, and suppose
\(\mathcal O_Y(D)\in\operatorname{Pic}^0(Y)\).  For every holomorphic vector
bundle \(B\), a general \(P\in\operatorname{Pic}^0(Y)\) satisfies
\[
 H^0(Y,B\otimes P)=0,\qquad
 H^{\dim Y}(Y,B\otimes P)=0.
\tag{5.23}
\]
\end{lemma}

\begin{proof}
The inclusions \(\mathcal O_Y(-(k+1)D)\hookrightarrow
\mathcal O_Y(-kD)\) give a descending chain
\[
 V_k=H^0(Y,B\otimes\mathcal O_Y(-kD))
 \subseteq H^0(Y,B).
\]
It stabilizes because the last vector space is finite-dimensional.  A
section in the stable value belongs to every \(V_k\).  In a local equation
\(z=0\) for \(D\), every coefficient of that section is divisible by
\(z^k\) for all \(k\), and hence is zero.  The section vanishes on a
neighbourhood of \(D\), and therefore on connected \(Y\) by analytic
continuation.  Thus \(V_k=0\) for some \(k\).  Since
\(\mathcal O_Y(-kD)\in\operatorname{Pic}^0(Y)\), upper semicontinuity in a
local universal family of topologically trivial line bundles shows that
\[
 \{P:h^0(Y,B\otimes P)>0\}
\]
is a proper analytic subset of \(\operatorname{Pic}^0(Y)\).  This proves
the first vanishing for general \(P\).  Serre duality identifies the dual
of the second group in (5.23) with
\[
 H^0(Y,B^\vee\otimes\omega_Y\otimes P^\vee).
\]
Apply the first part to \(B^\vee\otimes\omega_Y\) and use that inversion is
an automorphism of \(\operatorname{Pic}^0(Y)\).
\end{proof}

\begin{theorem}[corrected CDP global-vanishing criterion]
\label{thm:corrected-cdp-global}
\Verified\ Let \(X\) be a compact connected complex threefold such that
\[
 H^1(X;\mathbb Z)=H^2(X;\mathbb Z)=0,\qquad a(X)=1,
\]
and suppose its algebraic reduction is a holomorphic map
\(f:X\to C\).  If there exists \(L_0\in\operatorname{Pic}(X)\) with
\[
 H^2(X,T_X\otimes L_0)=0,
\tag{5.24}
\]
then, for \(M\) in a dense open subset of \(\operatorname{Pic}(X)\),
\[
 H^i(X,T_X\otimes M)=0\quad(i\ne1),
\tag{5.25}
\]
\[
 \chi(X,T_X\otimes M)\le0,\qquad
 \left\langle c_3(T_X),[X]\right\rangle\le0.
\tag{5.26}
\]
\end{theorem}

\begin{proof}
The exponential sequence and the two integral vanishings give
\[
 \operatorname{Pic}(X)=\operatorname{Pic}^0(X)
 \simeq H^1(X,\mathcal O_X).
\tag{5.27}
\]
Thus this Picard group is a connected irreducible complex vector group.  Any fibre divisor \(D\) of
\(f\) is nonzero and effective, and its line bundle lies in
\(\operatorname{Pic}^0(X)\) by \(H^2(X;\mathbb Z)=0\).  Applying
\cref{lem:cdp-generic-endpoints} to \(B=T_X\) shows that the degree-zero and
degree-three groups in (5.25) vanish away from two proper analytic subsets
of \(\operatorname{Pic}(X)\).

The cohomology-jump locus
\[
 \{M\in\operatorname{Pic}(X):h^2(X,T_X\otimes M)\ge1\}
\]
is analytic by upper semicontinuity.  It is proper because (5.24) gives a
point outside it.  Its complement therefore meets the two endpoint
vanishing loci in a dense open set.  This proves (5.25), and on that set
\[
 \chi(X,T_X\otimes M)=-h^1(X,T_X\otimes M)\le0.
\tag{5.28}
\]

Every \(M\in\operatorname{Pic}(X)\) has rationally trivial first Chern
class.  The hypotheses also give \(c_1(T_X)=0\) in rational cohomology and,
by Poincar\'e duality, \(H^4(X;\mathbb Q)=0\), so
\(c_2(T_X)=0\) rationally.  The degree-six term of
\(\operatorname{ch}(T_X\otimes M)\operatorname{td}(X)\) is consequently
\(\tfrac12c_3(T_X)\).  Hirzebruch--Riemann--Roch turns (5.28) into
\[
 \chi(X,T_X\otimes M)
 =\frac12\left\langle c_3(T_X),[X]\right\rangle\le0,
\]
which proves (5.26).
\end{proof}

\begin{corollary}[intact relative criterion]\label{cor:corrected-cdp-relative}
\Verified\ In the setting of
\cref{thm:corrected-cdp-global}, it is enough that one line bundle
\(L\in\operatorname{Pic}(X)\) satisfy
\[
 R^2f_*(T_X\otimes L)=0.
\tag{5.29}
\]
This is the valid content of the sufficient criterion printed as
\cite[Proposition~2.4, p.~680]{CDP2020}.
\end{corollary}

\begin{proof}
Let \(A\) be a positive-degree line bundle on the compact Riemann surface
\(C\), and put \(E=T_X\otimes L\).  Projection formula gives
\[
 R^2f_*(E\otimes f^*A^m)=R^2f_*E\otimes A^m=0.
\]
For \(m\gg0\), Serre vanishing on \(C\) gives
\[
 H^1(C,R^1f_*E\otimes A^m)=0.
\]
The Leray spectral sequence has only columns zero and one, so its
total-degree-two terms yield
\[
 H^2(X,E\otimes f^*A^m)=0.
\]
Thus \(L_0=L\otimes f^*A^m\) satisfies (5.24), and
\cref{thm:corrected-cdp-global} applies.
\end{proof}

\begin{theorem}[exact reduced-fibre repair of the p.~692 step]
\label{thm:corrected-cdp-fibrewise}
\Verified\ Retain the setting of
\cref{thm:corrected-cdp-global}, assume every full scheme-theoretic fibre
\(W_c=f^{-1}(c)\) is reduced, and fix \(L\in\operatorname{Pic}(X)\).  Put
\[
 A_c=(\omega_X\otimes L^\vee)|_{W_c},\qquad
 \mathscr T_c=\ker\!\left(
 \Omega^1_{W_c}\longrightarrow
 j_*\Omega^1_{(W_c)_{\rm reg}}\right).
\]
For each \(c\), let
\[
 \delta_c:
 H^0(W_c,\mathscr T_c\otimes A_c)
 \longrightarrow
 H^1(W_c,N^*_{W_c/X}\otimes A_c)
\tag{5.30}
\]
be the connecting map from the missing cotangent filtration.  If, for every
\(c\),
\[
 H^0(W_c,\widetilde\Omega^1_{W_c}\otimes A_c)=0,\qquad
 H^0(W_c,N^*_{W_c/X}\otimes A_c)=0,\qquad
 \delta_c\ \text{is injective},
\tag{5.31}
\]
then (5.29), (5.25), and (5.26) hold.
\end{theorem}

\begin{proof}
Apply \cref{prop:missing-global-section-criterion} to the whole reduced
Cartier fibre \(W_c\), not separately to its components.  Conditions (5.31)
give
\[
 H^0(W_c,\Omega_X^1|_{W_c}\otimes A_c)
 =\ker\delta_c=0.
\tag{5.32}
\]
The map \(f\) is flat: a local parameter at \(c\) pulls back to a
non-zero-divisor in every regular local ring of the connected manifold \(X\),
and torsion-freeness over the base DVR is flatness.  Moreover, the normal
line of the full fibre is the pullback of the one-dimensional vector space
\(T_cC\), hence is trivial.  Adjunction therefore identifies
\(\omega_{W_c}\simeq\omega_X|_{W_c}\).  With \(E=T_X\otimes L\), absolute
Cartier-fibre duality now gives
\[
 H^2(W_c,E|_{W_c})^\vee
 \simeq H^0(W_c,\Omega_X^1|_{W_c}\otimes A_c)=0.
\]
The uniform top-degree base-change isomorphism of
\cref{cor:surface-top-base-change} therefore gives
\[
 (R^2f_*E)_c\otimes_{\mathcal O_{C,c}}\mathbb C(c)=0
\]
for every \(c\).  Nakayama's lemma kills every stalk of the coherent sheaf
\(R^2f_*E\), proving (5.29).  The remaining conclusions follow from
\cref{cor:corrected-cdp-relative}.
\end{proof}

\begin{remark}[scope and sharpness of the repair]
Under the two endpoint vanishings in (5.31), injectivity of \(\delta_c\) is
equivalent to the desired ambient vanishing (5.32); it is not a politeness
assumption hiding the original claim.  For a reducible reduced fibre the
map must be computed on the whole Cartier fibre.  A nonreduced fibre requires
additional nilpotent-thickening layers and is not covered by
\cref{thm:corrected-cdp-fibrewise}.  The global criterion
\cref{thm:corrected-cdp-global} and the relative criterion
\cref{cor:corrected-cdp-relative} do allow nonreduced fibres.

The reconstructed candidate violates both added conditions rather than being
silently excluded by terminology.  Since \(c_3(X)=2\), the contrapositive of
\cref{thm:corrected-cdp-global} proves
\[
 H^2(X,T_X\otimes L_0)\ne0
 \quad\text{for every }L_0\in\operatorname{Pic}(X).
\]
Independently, \cref{cor:candidate-conductor-actual} proves
\[
 R^2f_*(T_X\otimes L)\ne0
 \quad\text{for every }L\in\operatorname{Pic}(X),
\]
and its cusp conductor supplies a nonzero element of the relevant
\(\ker\delta_c\).  The first conclusion rules out the global hypothesis; the
second independently rules out the stronger relative sufficient condition.
Thus the corrected theorem is compatible with, and
precisely separates, the verified counterexample to CDP20
Theorem~2.2 as stated.
\end{remark}

\paragraph{The exact monodromy repair.}
The split-extension argument can reach only the invariant integral lattice
\(H^1(F;\mathbb Z)^{\pi_1(C^\circ)}\).  Its final complex-linear inference is
valid precisely when
\[
 \operatorname{Span}_{\mathbb C}\operatorname{im}\!\left(
 H^1(F;\mathbb Z)^{\pi_1(C^\circ)}
 \longrightarrow H^1(F,\mathcal O_F)\right)
 =H^1(F,\mathcal O_F).
\tag{5.33}
\]
Trivial monodromy on \(H_1(F;\mathbb Z)\) is a stronger sufficient
hypothesis.  Neither (5.33) nor trivial monodromy follows from the hypotheses
of CDP20 Theorem~2.2.  The corrected theorem above bypasses this rank route
entirely; it does not relabel the false direct-sum formula as proved.

\subsection{Exact present disposition}

\begin{enumerate}[label=\arabic*.]
\item \Refuted\ The p.~692 reduction in \cite{CDP2020} is not a formal
consequence of the two cotangent sequences.
\Cref{prop:nodal-counterexample} is a proper geometric counterexample to
that exact sheaf-level inference for a non-torsion bundle on the central
fibre; no extension of that bundle to the displayed total space is claimed.
\item \Refuted\ The direct-product homology formula in
\cite[Lemma~3.2]{CDP1998} omits monodromy coinvariants.  This leaves the
torus branch imported by \cite{CDP2020} without its printed rank argument.
\item \Verified\ The normalization--conductor construction,
Cartier-fibre duality, flatness over the base curve, and top-degree analytic
base change are proved above under their stated hypotheses.
\item \Verified\ The period-family descent, toric and finite fillings,
global gluing, and integral analytic nearby-cycle/Leray bridges are proved in
\cref{sec:candidate-geometry,sec:integral-leray}.  Therefore the candidate's
hypotheses (H1)--(H4) hold and
$R^2f_*(T_X\otimes L)\ne0$ for every $L$ by
\cref{cor:candidate-conductor-actual}.
\item \Verified\ Those same construction and topology results, together
with \cref{thm:algebraic-dimension-audit}, hold for the constructed compact
complex threefold.  Hence
$H^1(X;\mathbb Z)=H^2(X;\mathbb Z)=0$, $a(X)=1$, $f$ is a holomorphic
algebraic reduction, and $c_3(X)=2$.  The resulting $X$ satisfies every
stated hypothesis of CDP Theorem~2.2 and contradicts its conclusion; it is
therefore a counterexample to that theorem as stated.  This theorem-level
conclusion uses the complete construction and topology proofs; it is not
inferred from the two printed proof defects alone.
\item \Verified\ The corrected global criterion
\cref{thm:corrected-cdp-global} bypasses both refuted dependencies:
existence of one \(L_0\) with
\(H^2(X,T_X\otimes L_0)=0\) implies generic off-degree-one vanishing and
\(c_3(X)\le0\).  The intact relative sufficient condition is
\(R^2f_*(T_X\otimes L)=0\), and
\cref{thm:corrected-cdp-fibrewise} replaces the invalid reduced-fibre step
by the exact injectivity condition on \(\delta_c\).  The original theorem
without these added hypotheses remains refuted by the constructed
threefold.  Independent specialist replication remains desirable as
external adjudication, but it is neither a mathematical hypothesis nor a
gate on the present proof program.
\end{enumerate}

\begin{test}[decisive theorem-level outcomes]\label{test:cdp-comparison}
The current workbench records two exact defects in the printed CDP chain, a
completed route from the constructed threefold to a theorem-level
counterexample, and the strongest corrected conclusion justified by the
intact argument.  The global vanishing in
\cref{thm:corrected-cdp-global} is false for every \(L_0\) by contraposition
from \(c_3(X)=2\).  Independently, the cusp-conductor calculation forces
\(R^2f_*(T_X\otimes L)\ne0\) for every \(L\), so the stronger relative
sufficient condition fails as well.  Thus the repair does not contradict the
counterexample or recover the original
theorem by renaming its conclusion.  Any further attempt to prove the
original statement must falsify a named step in the candidate reconstruction
or derive one of the new vanishing hypotheses from genuinely stronger
geometry.  Merely invoking the published theorem does neither.
\end{test}
